% **************************************************************************************************************
% A Classic Thesis Style
% An Homage to The Elements of Typographic Style
%
% Copyright (C) 2025 André Miede and Ivo Pletikosić
%
% If you like the style then I would appreciate a postcard. My address
% can be found in the file ClassicThesis.pdf. A collection of the
% postcards I received so far is available online at
% http://postcards.miede.de
%
% License:
% This program is free software; you can redistribute it and/or modify
% it under the terms of the GNU General Public License as published by
% the Free Software Foundation; either version 2 of the License, or
% (at your option) any later version.
%
% This program is distributed in the hope that it will be useful,
% but WITHOUT ANY WARRANTY; without even the implied warranty of
% MERCHANTABILITY or FITNESS FOR A PARTICULAR PURPOSE.  See the
% GNU General Public License for more details.
%
% You should have received a copy of the GNU General Public License
% along with this program; see the file COPYING.  If not, write to
% the Free Software Foundation, Inc., 59 Temple Place - Suite 330,
% Boston, MA 02111-1307, USA.
%
% PLEASE SEE ALSO THE AUTHORS' NOTE REGARDING THIS LICENSE
% IN THE DOCUMENTATION (ClassicThesis.pdf --> Chapter 1 / Chapter01.tex)
% **************************************************************************************************************
\RequirePackage{silence} % :-\
    \WarningFilter{scrreprt}{Usage of package `titlesec'}
    %\WarningFilter{scrreprt}{Activating an ugly workaround}
    \WarningFilter{titlesec}{Non standard sectioning command detected}
\documentclass[ twoside,openright,titlepage,numbers=noenddot,%1headlines,
                headinclude,footinclude,cleardoublepage=empty,abstract=on,
                BCOR=5mm,paper=a4,fontsize=11pt
                ]{scrreprt}

%********************************************************************
% Note: Make all your adjustments in here
%*******************************************************
\input{classicthesis-config}


%********************************************************************
% Bibliographies
%*******************************************************
\addbibresource{Bibliography.bib}
\addbibresource{papers/MRL/custom.bib}
\addbibresource[label=ownpubs]{ThesisWorks.bib}

%********************************************************************
% Hyphenation
%*******************************************************
%\hyphenation{put special hyphenation here}

% ********************************************************************
% GO!GO!GO! MOVE IT!
%*******************************************************
\begin{document}
\frenchspacing
\raggedbottom
\selectlanguage{american} % american ngerman
%\renewcommand*{\bibname}{new name}
%\setbibpreamble{}
\pagenumbering{roman}
\pagestyle{plain}
%********************************************************************
% Frontmatter
%*******************************************************
\include{FrontBackmatter/Titlepage}
%\include{FrontBackmatter/Titleback}
%\cleardoublepage\include{FrontBackmatter/Foreword}
\cleardoublepage\include{FrontBackmatter/Abstract}
\cleardoublepage%*******************************************************
% Works included in this thesis
%*******************************************************
\pdfbookmark[1]{Works included in this thesis}{publications}
\chapter*{Works included in this thesis}

% Dedicated sort keys preserve thesis order within the standard bibliography style.
\begin{refsection}[ownpubs]
    \nocite{thesis-lillama,thesis-baldwhisper,thesis-wolof,thesis-mrl,thesis-adaptive}
    \printbibliography[heading=none]
\end{refsection}

\cleardoublepage%*******************************************************
% Acknowledgments
%*******************************************************
\pdfbookmark[1]{Acknowledgments}{acknowledgments}

\begin{flushright}{\slshape
    Hammadi raconta au vieux mendiant \\
    le voyage extraordinaire qui le conduisit \\
    jusqu'auprès de Kaydara, le merveilleux Kaydara. \\
    Le vieux mendiant déclara à Hammadi : \\
    « Mon frère, apprends que chaque symbole recèle \\
    un sens, deux sens ou même plusieurs. \\
    Diurnes ou nocturnes peuvent être ces sens. \\
    Les diurnes sont fastes. \\
    Les nocturnes sont néfastes.\\
    Le caméléon que vous vîtes, [est le] premier signe \\
    chez les génies-nains qui servent au pays \\
    de Kaydara le merveilleux [...]} \\ 
    \medskip
    --- Hampâté Bâ, Amadou. (1969). \\
    \mbox{\textit{Kaïdara : récit initiatique peul}.}
\end{flushright}

\bigskip

\begingroup
\let\clearpage\relax
\let\cleardoublepage\relax
\let\cleardoublepage\relax
\chapter*{Acknowledgments}
\bigskip


\endgroup

\cleardoublepage\include{FrontBackmatter/Contents}
%********************************************************************
% Mainmatter
%*******************************************************
\cleardoublepage
\pagestyle{scrheadings}
\pagenumbering{arabic}
%\setcounter{page}{90}
% use \cleardoublepage here to avoid problems with pdfbookmark
\cleardoublepage
\ctparttext{%
\begin{flushright}{\slshape
    Le degré du caméléon est appelé vestibule royal; \\
    en effet tout s'y rencontre. \\
    Certains viennent donner le mal et son espèce. \\
    Certains viennent quémander, d'autres sont là pour offrir. \\
    Certains viennent mentir, d'autres parce qu'on a menti sur eux.} \\ \medskip
    --- Hampâté Bâ, Amadou. (1969). \\
    \mbox{\textit{Kaïdara : récit initiatique peul}.}
\end{flushright}}
\part{Introduction}\label{pt:manual}
%************************************************
\chapter{Modeling Language}\label{ch:introduction}
%************************************************
\section{A dialog between Shannon and Chomsky}
It is curious that we still don't know where the expression \textit{language model} comes from or when it became established.
We mean the \textit{expression} itself, rather than the concept.
The concept of a language model is quite old, and one of the earliest
formulations closely related to what we now call language modeling can be found in Shannon's 
work on information theory~\citep{shannon1948mathematical}\footnote{
    This intuition was already explored by Markov in 1913~\citep{1913_lm};
    Shannon later provided a more fully developed formalization.
}. 
Shannon starts from a simple observation: sequences of linguistic symbols are not random, 
but exhibit regularities characteristic of the language.
\begin{quote}
    \small\itshape
    ``[\ldots{}] the messages to be transmitted consist of sequences of letters. 
    These sequences, however, are not completely random. 
    In general, they form sentences and have the statistical structure of, say, English.''

    \begin{flushright}
    \normalfont--- Claude E. Shannon, \textit{A Mathematical Theory of Communication (1948)}
    \end{flushright}
\end{quote}
If English sentences have a statistical structure, then a statistical model that 
captures this structure should be able to generate sequences that look like English.
But what is this statistical model? How can we estimate it? Later, Shannon gives a more operational view:

\begin{quote}
    \small\itshape
    ``We can think of a discrete source as generating the message, symbol by symbol.
    It will choose successive symbols according to certain probabilities depending,
    in general, on preceding choices [\ldots{}].''

    \begin{flushright}
    \normalfont--- Claude E. Shannon, \textit{A Mathematical Theory of Communication (1948)}
    \end{flushright}
\end{quote}
Shannon thus describes a statistical
model of English but does not call it a \textit{language model}.
The story takes a different turn in the early 1970s. In 1974, a team at IBM working 
on speech recognition described a statistical model closely related to Shannon's formulation and explicitly 
called it a \textit{language model}. In their probabilistic decomposition of speech recognition, 
the language model is simply the component that assigns probabilities to possible 
word sequences---that is, a model of which utterances are more or less likely in a given language. 
As they put it, ``it is the function of the language model (LM) to provide us with $\Pr\{W \mid L\}$,'' 
where $W$ denotes a word sequence and $L$ the language. 
This work is therefore probably one of the earliest uses of the expression \textit{language model} 
in essentially its modern statistical sense: a probability distribution over linguistic sequences 
estimated from language data.

Calling these systems language models could suggest that they are
linguistic models of a language. But does a better statistical model of English mean a better model of the language itself?
The idea that statistical learning can capture linguistic structure is compatible 
with usage-based linguistics and connectionist approaches. Nativist approaches, 
however, question whether statistical learning alone is sufficient to explain linguistic knowledge.
In \textit{Three Models for the Description of Language}, Chomsky states~\citep[p.~113]{chomsky1956three}:


\begin{quote}
    \small\itshape
    ``We find that no finite-state Markov process that produces symbols with transition from state to state can serve as an English grammar. 
    [...] such processes that produce $n$-order statistical approximations to English do not come closer, with increasing $n$, to matching the output of an English grammar.''

    \begin{flushright}
    \normalfont--- Noam Chomsky, \textit{Three Models for the Description of Language (1956)}
    \end{flushright}
\end{quote}

There are two main claims made by Chomsky in this paper. 
The first claim is structural and concerns the \textbf{bounded memory} of finite-state models.
For finite-order Markov language models, this limitation manifests as a fixed context length: 
only a bounded portion of the preceding sequence can determine the next symbol.
English, however, permits recursive constructions whose dependencies can require information 
from an arbitrarily large preceding structure.
Modern language models have greatly increased the amount of preceding context they can exploit, 
making them much better at modeling long-range dependencies. 
Yet practical models still operate with bounded context and memory, 
so the distinction highlighted by Chomsky between 
unbounded linguistic structure and bounded computational memory remains relevant. 
Second, even if we consider a statistical model with unbounded memory, 
Chomsky's criticism goes beyond model capacity. His broader claim is that \textbf{language is innate},
biologically programmed in the human brain~\citep{chomsky1965aspects} and 
cannot be learned simply by accumulating statistics over observed linguistic data.
The main argument is the \textit{poverty of stimulus}: a child learns their native language rapidly,
despite being exposed to limited, noisy, and incomplete examples. According to this account, infants have structural biases that help them learn and generalize linguistic structures from a few examples.
Chomsky's theory holds that human linguistic competence depends on internal constraints 
that cannot be inferred from experience alone.

It is not straightforward, however, to translate this notion of innateness into 
the language of machine learning. Does Chomsky's argument imply that general-purpose 
learning architectures are inherently incapable of acquiring language from data alone? 
Or can we instead design computational architectures and general learning objectives whose inductive biases make
the acquisition of linguistic structure easier?
\footnote{For Chomsky, this question is not really interesting, 
    as reproducing linguistic behavior with an artificial architecture is not, 
    by itself, an explanation of human language acquisition: 
    \url{https://www.youtube.com/watch?v=PZQOIuuQiJA}}
Maybe we can also reverse the question: can learning from experience alone help
uncover the structure latent in language? Even though modern machine learning algorithms are still far from matching
the data efficiency of human children, substantial progress has been made by
designing learning architectures that optimize the language modeling objective, improving
the linguistic competence of these systems.

The language modeling objective, as Shannon describes it, is to learn to predict
the next symbol from a fixed number of previous symbols.
These symbols can be letters or words, but nowadays they are usually tokens, units that can form a word or a subword. 
More formally, the probability of picking a token $s_i$ given a sequence of previous tokens $s_{<i} = (s_1, s_2, \ldots, s_{i-1})$ is written as $P(s_i \mid s_{<i})$.
Using the chain rule, these individual probabilities can be multiplied together to give the probability of the whole text:
\begin{equation}
    P(s) = p(s_1) \times p(s_2 | s_1) \times p(s_3 | s_{1}, s_{2}) \times \cdots \times p(s_t | s_{<t}) = \prod_{i=1}^t P(s_i \mid s_{<i}).
\end{equation}
Each individual probability can be defined from a score over the vocabulary:
\begin{equation}
    p(s_i \mid s_{<i}) =
    \frac{\mathrm{score}(s_{<i}, s_i)}
    {\sum_{v \in \mathcal{V}} \mathrm{score}(s_{<i}, v)}.
\end{equation}
Here, $\mathrm{score}_\theta(x, y)$ measures how strongly the context $x$ is associated with the next token $y$, and $\mathcal{V}$ is the token vocabulary.
The function $\mathrm{score}(x, y)$ can be designed very freely but in neural language models, it is often the exponential of the dot product
between the context vector of $x$ and the token vector $y$.

Using maximum likelihood estimation, we can estimate sequence probabilities with a statistical model parameterized
by $\theta$ and trained on a large corpus $\mathcal{C}$ of texts:
\begin{equation}
    \theta^* = \arg\max_\theta
    \sum_{s \in \mathcal{C}} \sum_{i=1}^{|s|}
    \log P_\theta(s_i \mid s_{<i}).
\end{equation}

This is the learning objective of language modeling; it remains to design a learning algorithm 
that can optimize it effectively. While many neural architectures have been used for this 
purpose, ranging from CNNs to RNNs, one architecture has proven particularly effective: 
the Transformer. We next present the Transformer through the lens of linguistic 
and knowledge acquisition, and discuss how scaling this architecture leads to further gains 
in knowledge and linguistic competence.

\section{The Transformer, the best machine language learner?}
Earlier work improved neural machine translation by adding attention
to an RNN-based encoder--decoder~\citep{bahdanau2015neural}. The Transformer
replaced recurrence with attention~\citep{vaswani2017attention}. Removing recurrent dependencies allows all token positions
within a layer to be processed in parallel during training, making the architecture more
parallelizable and easier to scale. Figure~\ref{fig:decoder-only-transformer} presents the
decoder-only Transformer architecture used for autoregressive language modeling.
The model stacks $N$ layers, each containing a self-attention module followed by an MLP,
before an output head computes the next-token probability distribution. We next discuss
these three components and their roles in acquiring knowledge and generalizing linguistic
structures.

\begin{figure}[ht]
    \centering
    \includegraphics[width=0.75\linewidth]{figures/transformer.pdf}
    \caption{Decoder-only Transformer architecture.}
    \label{fig:decoder-only-transformer}
\end{figure}

\subsection{Attention}
The innovation of the Transformer is to remove the RNN and replace it entirely with
\textbf{Multi-Head Attention (MHA)}. Given $\mathbf{X} \in \mathbb{R}^{T \times d_{\mathrm{model}}}$,
the sequence of token embeddings for $T$ positions, MHA projects the sequence into queries, keys
and values for $H$ heads:
\begin{equation}
    \mathbf{Q}, \mathbf{K}, \mathbf{V}
    \in \mathbb{R}^{T \times H \times d_{\mathrm{head}}}.
\end{equation}
Equivalently, each head $h$ has its own projections:
\begin{align}
    \mathbf{Q}_h &= \mathbf{X}\mathbf{W}^{Q}_h, &
    \mathbf{K}_h &= \mathbf{X}\mathbf{W}^{K}_h, &
    \mathbf{V}_h &= \mathbf{X}\mathbf{W}^{V}_h,
\end{align}
where $\mathbf{W}^{Q}_h \in \mathbb{R}^{d_{\mathrm{model}} \times d_q}$, 
$\mathbf{W}^{K}_h \in \mathbb{R}^{d_{\mathrm{model}} \times d_k}$
and $\mathbf{W}^{V}_h \in \mathbb{R}^{d_{\mathrm{model}} \times d_v}$ are learned projection
matrices for head $h$. The head output is then computed with scaled dot-product self-attention:
\begin{equation}
    \mathrm{head}_h
    =
    \mathrm{softmax}\left(
        \frac{\mathbf{Q}_h \mathbf{K}_h^\top + \mathbf{M}}{\sqrt{d_k}}
    \right)\mathbf{V}_h .
\end{equation}
For a decoder language model, $\mathbf{M}$ is a causal mask: entries corresponding to future
positions are set to $-\infty$, so token $i$ can only attend to tokens $j \leq i$. Multi-head attention runs $H$ such heads in parallel and concatenates their outputs:
\begin{equation}
    \mathrm{MHA}(\mathbf{X})
    =
    \mathrm{Concat}\left(\mathrm{head}_1,\ldots,\mathrm{head}_H\right)\mathbf{W}^{O},
\end{equation}
with $\mathbf{W}^{O} \in \mathbb{R}^{H d_v \times d_{\mathrm{model}}}$. Each head can therefore
learn a different projection of the same context.

Having multiple independent heads is beneficial since it has been shown
that \textbf{different attention heads play linguistically interpretable roles}.
Some specialize in positional relations, syntactic
dependencies, coreference, or other linguistically interpretable patterns
\citep{clark-etal-2019-bert,voita-etal-2019-analyzing}. However, attention
heads are also substantially redundant, as many of them can be removed with
little or no degradation in model performance
\citep{michel-etal-2019-sixteen,voita-etal-2019-analyzing,
kovaleva-etal-2019-revealing}.


In standard MHA, as introduced in the original Transformer paper, the projection matrices
$\mathbf{W}^{Q}_h$, $\mathbf{W}^{K}_h$ and $\mathbf{W}^{V}_h$ are usually chosen so that all
heads together project back to a space of size $d_{\mathrm{model}}$. However, these projections
can become a bottleneck during generation. In a naive implementation, at each decoding step, all
previous tokens would be projected again to recompute their keys and values. The
\textbf{KV-Cache} reduces this computation and speeds up generation by caching the keys and values
of previous tokens. However, this cache grows linearly with the sequence length, and it is easy to
run out of memory when generating very long texts. To reduce the KV-Cache, we need to reduce the
dimensionality of keys and values.
Among the popular approaches to reduce the KV-Cache, \textbf{Grouped-Query Attention (GQA)}
shares key-value pairs across multiple query heads. Intuitively, the keys and values have fewer heads
than the queries:
\begin{align}
    \mathbf{K}, \mathbf{V} &\in \mathbb{R}^{T \times h_g \times h_{\mathrm{dim}}}, \\
    \mathbf{Q} &\in \mathbb{R}^{T \times h_q \times h_{\mathrm{dim}}},
    \qquad h_g < h_q .
\end{align}
The key-value heads are then broadcast to the query heads, and attention is computed as usual:
\begin{equation}
    \mathrm{Attention}(\mathbf{Q}, \mathbf{K}, \mathbf{V})
    =
    \mathrm{softmax}\left(\mathbf{Q}\mathbf{K}^{\top}\right)\mathbf{V}.
\end{equation}
Standard MHA corresponds to no sharing, while multi-query attention is the extreme case where all
query heads share the same keys and values. GQA is the compromise between the two: it reduces the
KV-Cache while keeping several different key-value heads.
While GQA reduces KV-cache memory, it does not reduce the attention computation. 
Nowadays, GQA is often combined with \textbf{Sparse Attention}, 
such as DeepSeek Sparse Attention~\citep{deepseekai2025deepseekv32pushingfrontieropen}, to reduce this compute 
by allowing each query to attend to only a subset of previous tokens.


\subsection{Multilayer Perceptron}
After the attention layer, the contextualized tokens are fed into the multilayer perceptron (MLP). While attention mixes information across tokens, the MLP mixes information across channels: each token is processed independently. In the original Transformer paper, the MLP is defined as:
$$
\text{MLP}(x) = f(\textbf{x}\textbf{W}_{up})\textbf{W}_{down}
$$
where $f(\cdot)$ is the non-linear activation function, 
$\textbf{W}_{\text{up}} \in \mathbb{R}^{d_{\text{model}} \times d_{\text{inter}}}$ 
projects the input token $\mathbf{x} \in \mathbb{R}^{d_{\text{model}}}$ into a higher-dimensional space,
with $d_{\text{model}} \ll d_{\text{inter}}$. Finally, $\textbf{W}_{\text{down}} \in \mathbb{R}^{d_{\text{inter}} \times d_{\text{model}}}$,
projects back to the model dimension.
Modern LLMs often replace this simple MLP with a gated linear unit variant. Instead of applying
only one non-linearity to the expanded representation, the block learns two projections: one
projection $\mathbf{V}_{\text{up}}$ produces the values, and the other projection
$\mathbf{W}_{\text{up}}$ gates them. Omitting bias terms, these alternatives can be written as:
\begin{align}
    \mathrm{MLP}_{\mathrm{GLU}}(\mathbf{x})
    &= \left(\sigma(\mathbf{x}\mathbf{W}_{\text{up}}) \odot \mathbf{x}\mathbf{V}_{\text{up}}\right)\mathbf{W}_{\text{down}}, \\
    \mathrm{MLP}_{\mathrm{GEGLU}}(\mathbf{x})
    &= \left(\mathrm{GELU}(\mathbf{x}\mathbf{W}_{\text{up}}) \odot \mathbf{x}\mathbf{V}_{\text{up}}\right)\mathbf{W}_{\text{down}}, \\
    \mathrm{MLP}_{\mathrm{Bilinear}}(\mathbf{x})
    &= \left(\mathbf{x}\mathbf{W}_{\text{up}} \odot \mathbf{x}\mathbf{V}_{\text{up}}\right)\mathbf{W}_{\text{down}}, \\
    \mathrm{MLP}_{\mathrm{SiLU}}(\mathbf{x})
    &= \left(\mathrm{SiLU}(\mathbf{x}\mathbf{W}_{\text{up}}) \odot \mathbf{x}\mathbf{V}_{\text{up}}\right)\mathbf{W}_{\text{down}}, \\
    \mathrm{MLP}_{\mathrm{ReGLU}}(\mathbf{x})
    &= \left(\mathrm{ReLU}(\mathbf{x}\mathbf{W}_{\text{up}}) \odot \mathbf{x}\mathbf{V}_{\text{up}}\right)\mathbf{W}_{\text{down}}.
\end{align}
Here, $\odot$ denotes element-wise multiplication. Table~\ref{tab:gated_mlp} gives the performance of each activation on the SuperGLUE benchmark.
The SiLU-based variant gives better results, which helps explain why many LLMs, including Mistral-v0.1 and Granite, use it.
\begin{table*}[h!]
    \centering
    \setlength{\tabcolsep}{8pt}
    \begin{tabular}{>{\raggedright\arraybackslash}p{0.30\textwidth}>{\centering\arraybackslash}p{0.40\textwidth}}
        \toprule
        \textbf{} & \textbf{Average SuperGLUE score} \\
        $\mathrm{MLP}_{\mathrm{GLU}}$ & 73.95 \\
        $\mathrm{MLP}_{\mathrm{GEGLU}}$ & 73.96 \\
        $\mathrm{MLP}_{\mathrm{Bilinear}}$ & 73.81 \\
        \rowcolor[HTML]{EEE0CC}$\mathrm{MLP}_{\mathrm{SwiGLU}}$ & \textbf{74.56} \\
        $\mathrm{MLP}_{\mathrm{ReGLU}}$ & 73.66 \\
        \bottomrule
    \end{tabular}
    \caption{Average SuperGLUE score for gated MLP variants.}
    \label{tab:gated_mlp}
\end{table*}
It is difficult to explain why gated MLPs work better with some activations than others.
As the author of the study puts it~\citep{shazeer2020gluvariants}:
\begin{quote}
    \small\itshape
    ``We offer no explanation as to why these architectures seem to work; we attribute their
    success, as all else, to divine benevolence.''

    \begin{flushright}
    \normalfont--- Noam Shazeer, \textit{GLU Variants Improve Transformer (2020)}
    \end{flushright}
\end{quote}
In modern LLM architectures, MLP parameters can account for over $90\%$ of the total model parameters.
It has been shown that MLP layers play an important role in storing knowledge acquired during web-scale pretraining. 
One interpretation treats the MLP as a key--value memory~\citep{geva-etal-2021-transformer}:
each column of \(\mathbf{W}_{\text{up}}\) in the MLP represents a key
\(\mathbf{k}_i\) that detects linguistic patterns seen during training, while
each row of \(\mathbf{W}_{\text{down}}\) represents the corresponding value
\(\mathbf{v}_i\), which promotes the tokens likely to follow those patterns.
The key \(\mathbf{k}_i\) can be strongly or weakly activated depending on the
input \(x\). When \(\mathbf{k}_i\) is strongly activated, its corresponding
value \(\mathbf{v}_i\) is assigned a large weight and therefore contributes
more strongly to the MLP output. This mechanism can be schematized as follows:
\[
\boxed{
\underbrace{\mathbf{k}_i(\mathbf{x})}_{
    \text{whether the key's pattern is active}
}
\;\longrightarrow\;
\underbrace{\mathbf{v}_i}_{
    \text{promotes tokens likely to follow the pattern}
}
}
\]
\begin{table*}[t]
    \centering
    \small
    \setlength{\tabcolsep}{6pt}
    \renewcommand{\arraystretch}{1.15}

    \begin{tabularx}{\textwidth}{@{} l l c >{\raggedright\arraybackslash}X @{}}
        \toprule
        \textbf{Value}
        & \textbf{Top prediction}
        & \textbf{Precision@50}
        & \textbf{Matching trigger example} \\
        \midrule

        $\mathbf{v}^{15}_{222}$
        & \emph{each}
        & 68\%
        & \emph{But when bees and wasps resemble \textcolor{blue}{each}} \\

        \midrule

        $\mathbf{v}^{16}_{752}$
        & \emph{played}
        & 16\%
        & \emph{Her first role was in Vijay Lalwani's psychological thriller
        Karthik Calling Karthik, where Padukone was cast as the supportive
        girlfriend of a depressed man (\textcolor{blue}{played}} \\

        \midrule

        $\mathbf{v}^{13}_{2601}$
        & \emph{extratropical}
        & 4\%
        & \emph{Most of the winter precipitation is the result of synoptic-scale,
        low-pressure weather systems (large-scale storms such as
        \textcolor{blue}{extratropical}} \\

        \midrule

        $\mathbf{v}^{15}_{881}$
        & \emph{part}
        & 92\%
        & \emph{Comet served only briefly with the fleet, owing in large
        \textcolor{blue}{part}} \\

        \midrule

        $\mathbf{v}^{16}_{2070}$
        & \emph{line}
        & 84\%
        & \emph{Sailing from Lorient in October 1805 with one ship of the
        \textcolor{blue}{line}} \\

        \midrule

        $\mathbf{v}^{12}_{3186}$
        & \emph{jail}
        & 4\%
        & \emph{On May 11, 2011, four days after scoring 6 touchdowns for the
        Slaughter, Grady was sentenced to twenty days in
        \textcolor{blue}{jail}} \\

        \bottomrule
    \end{tabularx}

    \caption{
    Examples of MLP key--value memory pairs identified in prior work
    \citep{geva-etal-2021-transformer}.
    For each value vector $\mathbf{v}^{\ell}_{i}$, \emph{Top prediction}
    is the vocabulary token receiving the highest score when the value is
    projected through the model's output embedding.
    \emph{Precision@50} is the percentage of the 50 training prefixes that
    most strongly activate the paired key $\mathbf{k}^{\ell}_{i}$ whose
    ground-truth next token equals that prediction.
    The final column presents one such matching prefix, with its
    ground-truth next token shown in blue.
    }
    \label{tab:ffn-key-value-examples}
\end{table*}
The MLP does not only store lexical knowledge. In some cases, this lexical knowledge coincides
with factual knowledge. When a key detecting the prefix \textit{The capital of France is} 
is strongly associated with a value representing \textit{Paris}, this lexical key--value memory 
turns out to also encode world knowledge. 
Hence, more broadly, it is widely believed that MLPs play a central role in storing the knowledge acquired during pretraining
\citep{geva-etal-2021-transformer,dai-etal-2022-knowledge}.
This view is also reflected in many knowledge-editing methods, which often directly target MLP weights
\citep{meng2023locatingeditingfactualassociations,meng2023masseditingmemorytransformer, fang2025alphaeditnullspaceconstrainedknowledge}.
If we view the MLP as a key--value memory, it is reasonable to expect that each input token 
requires access to only a subset of its key--value pairs. We will later discuss how this 
intuition can guide the design of an adaptive sparse memory layer.

\subsubsection{Language Modeling Head}
Well, the tokens are contextualized thanks to the multi-head attention, and each token retrieves 
the relevant knowledge encoded within the MLP. This process is repeated for $L$ layers and at the end
produces a context vector $\textbf{c}_{<i}$ that is ready for next-token prediction
using the language modeling head (LM Head or output embedding):

\begin{equation}
    p(s_i \mid s_{<i}) =
    \frac{\mathrm{\exp}(\mathbf{c}_{<i} \mathbf{W}_i)}
    {\sum_{v \in \mathcal{V}} \mathrm{\exp}(\mathbf{c}_{<i} \mathbf{W}_v)}.
\end{equation}
where $W \in \mathbb{R}^{d \times |\mathcal{V}|}$ is the LM Head weight matrix.
The LM head weight matrix is another important component of large-scale Transformer LLMs. It is closely linked to the tokenizer:
the vocabulary size determines the number of parameters in the output projection, 
while the quality of tokenization affects how easily the model learns and, ultimately, its performance.
Recent tokenizers use large vocabularies to capture the diversity of tokens in web-scale data.
Hence, the LM head can account for a large share of the parameters in an LLM, particularly in small models. Figure \ref{fig:emb-param-ratio} shows the memory
occupied by the input/output embedding parameters for different model sizes, including dense and quantized ones.
We can see that the embedding and LM Head parameters represent a larger fraction of the model 
when the Transformer backbone is small\footnote{Quantization is usually not applied to the embedding weights.}.
\begin{figure}[ht]
    \centering
    \includegraphics[width=1.0\linewidth]{figures/embedding_parameter_ratio.pdf}
    \caption{Share of model storage occupied by embedding parameters across different model sizes.
    Data taken from \href{https://blog.squeezebits.com/vocabulary-trimming-methods}{https://blog.squeezebits.com/vocabulary-trimming-methods}.}
    \label{fig:emb-param-ratio}
\end{figure}
This is mainly caused by the large vocabulary of the models. A common solution to reduce the size by half is
weight tying \citep{press-wolf-2017-using}, which shares the input and output embedding matrices. Almost all small LLMs use weight tying, which can substantially
reduce the model's parameter count without hurting performance.

Another approach to reducing the number of parameters occupied by the LM head
is to factorize its weight matrix into two low-rank matrices. This strategy was
notably adopted in ALBERT~\citep{lan2020albertlitebertselfsupervised}, where
the embedding dimension was reduced from $768$ to $128$ without sacrificing downstream performance. 
More precisely, given the input $\mathbf{x}\in\mathbb{R}^{d}$ and an LM-head weight
matrix $\mathbf{W}\in\mathbb{R}^{d\times V}$, the factorized projection is
\begin{equation}
    \mathbf{x}\mathbf{W}
    \quad\longrightarrow\quad
    (\mathbf{x}\mathbf{W}_1)\mathbf{W}_2,
\end{equation}
where $\mathbf{W}_1\in\mathbb{R}^{d\times r}$,
$\mathbf{W}_2\in\mathbb{R}^{r\times V}$, and $r\ll H$. This reduces the
LM-head parameter count from $dV$ to $dr+rV$.

However, this bottleneck architecture can hurt optimization.
Some works have identified what is known as the
\textit{softmax bottleneck}~\citep{yang2018breakingsoftmaxbottleneck}:
a low-dimensional LM head can only produce a limited variety of output
distributions. Therefore, reducing its rank too aggressively can limit the
model's ability to predict complex token distributions. A study of LM heads
with different ranks found that reducing the rank
hurts performance and limits the expressiveness of the output distribution~\citep{godey2024smalllanguagemodelsunderperform}, as shown in Figure~\ref{fig:lm_head_botlnck}.
More importantly, this also hurts the learning signal during training \citep{godey2026lostbackpropagationlmhead}.
There are many solutions to avoid the \textit{softmax bottleneck}. Since this phenomenon arises mainly when
$r \ll V$, an obvious solution is to increase $r$ or reduce the vocabulary size $V$. Reducing the vocabulary size can be
achieved using a tokenizer with fewer tokens or falling back to a character- or byte-level tokenizer.
It is also technically possible after training to remove tokens from the input and output embeddings.
But this is risky as it can affect critical tokens. This discussion is related to the research we will present
later, which proposes to reduce the rank of the LM Head instead of pruning tokens in multilingual transformer models.
\begin{figure}[ht]
    \centering
    \subfloat[Accuracy.]{
        \includegraphics[width=0.48\linewidth]{figures/figure6_accuracy.png}
    }
    \hfill
    \subfloat[Cross-entropy.]{
        \includegraphics[width=0.48\linewidth]{figures/figure6_cross_entropy.png}
    }
    \caption{Effect of reducing the LM Head rank on accuracy and cross-entropy.
    Figures from \citep{godey2024smalllanguagemodelsunderperform}.}
    \label{fig:lm_head_botlnck}
\end{figure}

\subsection{A Transformer Is Never Large Enough}
What made the Transformer architecture particularly successful in language
modeling is its ability to scale both the number of parameters and the amount
of training data. Since Transformer training can be efficiently parallelized
across multiple GPUs, increasingly large models can be trained on increasingly
large datasets, leading to improved performance and the emergence of new
capabilities. However, scaling cannot be done blindly: given the high cost of
pretraining, the allocation between model size and training data must be
estimated before the final training run. For a fixed compute budget, how large
should the model be, and on how many tokens should it be trained?

Denoting the number of parameters by $N$, the number of training tokens by $D$,
and the training compute by $C$, the compute cost of a dense Transformer can be
approximated as $C \approx 6ND$. Given a fixed training budget $C$
(the number of floating-point operations available for training), how should
we allocate compute between model size and training tokens?
Both Kaplan \citep{kaplan2020scaling} and Chinchilla \citep{hoffmann2022training} 
scaling laws indicate that the compute-optimal model size and dataset size follow power laws
in the compute budget: $N_{\mathrm{opt}} \propto C^{\alpha}$ and $D_{\mathrm{opt}} \propto C^{\beta}$.
\textbf{Kaplan scaling laws} estimate the compute-optimal allocation as
$N_{\mathrm{opt}} \propto C^{0.73}$ and  $D_{\mathrm{opt}} \propto C^{0.27}$.
Thus, model size should increase much faster than training data, 
which favors very large models trained on comparatively few tokens. 
On the other hand, the \textbf{Chinchilla scaling laws}
later found that model size and training data should instead scale approximately equally:
$N_{\mathrm{opt}} \propto C^{0.5}$ and $D_{\mathrm{opt}} \propto C^{0.5}$.

Beyond guiding compute allocation, scaling laws also reveal how model capabilities
evolve with scale. While performance often improves predictably as training compute
increases, some abilities appear only after the model reaches a certain scale. The
model may remain close to random performance across several scales before suddenly
crossing a plateau and acquiring a new ability. Several such
\textit{emergent abilities} have been identified~\citep{wei2022emergent}, as shown in Figure~\ref{fig:wei-emergent-abilities}.
\begin{figure}[ht]
    \centering
    \includegraphics[width=1.0\linewidth]{figures/wei_emergent_abilities_figure.pdf}
    \caption{Examples of emergent abilities in large language models.
    Figure from \citep{wei2022emergent}. 
    An ability is emergent when its score surpasses the random baseline
    after remaining near that baseline at smaller scales.}
    \label{fig:wei-emergent-abilities}
\end{figure}
We can see that many linguistic skills such as language understanding, especially
for low-resource languages, only show up at a given model scale.
Then a question arises. Is there a point beyond which increasing the number of parameters no longer
produces meaningful improvements? Testing this hypothesis directly is difficult,
since training models at the largest scales is extremely costly.
Classical scaling laws predict \textit{diminishing returns}: test loss continues
to decrease with scale, but each additional increase in model size produces a
smaller improvement~\citep{kaplan2020scaling,hoffmann2022training}. However,
whether this implies diminishing returns in model capabilities depends on how
performance is measured. It has been argued that short-task
benchmarks may create an illusion of saturation~\citep{sinha2026illusion}. Even small improvements in
single-step accuracy can compound into large improvements in the length of tasks
that a model can reliably execute. Therefore, diminishing returns on loss or
short-task accuracy do not necessarily imply diminishing returns in practical
capabilities.
\include{Chapters/Chapter02_EfficientLanguageModels}
\cleardoublepage
\ctparttext{%
\begin{flushright}{\slshape
    Ce que signifient quelques traits du caméléon \\
    en diurne et nocturne, Hammadi, le voici. \\
    Changer de couleur, signifie en diurne \\
    qu'on est un homme sociable et agréable \\
    plein de tact, patient, commode à vivre \\
    et capable de s'adapter à toute circonstance, \\
    de l'accepter telle quelle et d'où qu'elle vienne ; \\
    être sociable et agréer les us et coutumes d'autrui. \\
    Dans le sens nocturne, cela devient hypocrisie, \\
    manque de constance et changement instable \\
    au point d'être impossible avec ses compagnons \\
    et de n'avoir ni personnalité ni dignité.} \\ \medskip
    --- Hampâté Bâ, Amadou. (1969). \\
    \mbox{\textit{Kaïdara : récit initiatique peul}.}
\end{flushright}}
\part{Compression of Transformer Models}\label{pt:compression}
%*****************************************
\chapter{Lillama: Large Language Models Compression via Low-Rank Feature Distillation}\label{ch:lillama}
%*****************************************

\section{Introduction}

\begin{figure}[ht]
        \centering
        \includegraphics[width=\linewidth]{papers/lillama/overview.pdf}
        \caption{\textbf{Lillama approach}: \textbf{STEP 1} selects layers to compress for a target compression ratio (e.g., N\%) using various strategies (see Section \ref{sec:algo}). \textbf{STEP 2} compresses and initializes the chosen parameters via SVD. \textbf{STEP 3} distills the low-rank weights with a small calibration dataset.}
        \label{fig:overview}
\end{figure}


\begin{table*}[h!]
\setlength{\tabcolsep}{2pt}
\begin{tabular}{>{\raggedright\arraybackslash}p{0.32\textwidth}|>{\centering\arraybackslash}p{0.12\textwidth}|>{\centering\arraybackslash}p{0.12\textwidth}|>{\centering\arraybackslash}p{0.12\textwidth}|>{\centering\arraybackslash}p{0.12\textwidth}|>{\centering\arraybackslash}p{0.12\textwidth}}
\toprule
\textbf{Method} & \textbf{No Custom Kernel} & \textbf{No Ampere-only} & \textbf{Memory Gain} & \textbf{Is One-Shot} & \textbf{Calib. Data} \\  \hline
SparseGPT \citep{frantar2023sparsegptmassivelanguagemodels} & \xmark & \xmark & \xmark & \cmark & - \\
Wanda \citep{wanda} & \xmark & \xmark & \xmark & \cmark & - \\
Teal \citep{TEAL} & \xmark & \xmark & \xmark & \cmark & - \\
ShearedLLaMA \citep{xia2024sheared} & \cmark & \cmark & \cmark & \xmark & 52B\\ 
Minitron \citep{minitron} & \cmark & \cmark & \cmark & \xmark & 94B \\
\rowcolor[HTML]{EEE0CC}\textbf{Lillama} & \cmark & \cmark & \cmark & \cmark & \textbf{0.013B} \\
\bottomrule
\end{tabular}
\caption{Positioning of our approach within the compression literature. \textbf{No Custom Kernel} indicates that a custom GPU kernel is not required to observe speedup and \textbf{No Ampere-only} indicates that speedup is not restricted to Ampere GPUs. \textbf{Calib. Data} is the total number of tokens in billions used for compression.}
\label{tab:litt}
\end{table*}


As discussed earlier, compressing large language models is a crucial challenge that can be addressed through several complementary approaches: quantization, pruning, and distillation. In this section, we introduce a new one-shot compression method that locally distills low-rank weights and satisfies the following three objectives:
\begin{enumerate}
\item The compression algorithm must be compute-efficient, e.g., it should require minimal computational resources and run in a reasonable amount of time. For instance, iterative magnitude pruning as used in prior work~\citep{frankle2019lottery} is too costly to be applied to large-scale LLMs.
\item It must achieve fast convergence using as little training data as possible; this objective is difficult to meet with many model pruning methods that require significant retraining after pruning to recover the lost performance.
\item The method should not cause a significant drop in performance; ideally, all the capabilities of the original LLM should remain intact.
\end{enumerate}
To address the first constraint, we compress the model with a local distillation objective, as opposed to a global objective that is more costly. We further reduce the cost by compressing and distilling only a subset of all layers.
To address the second constraint, we initialize the compressed layers with Singular Value Decomposition (SVD), which enables
reducing the required number of gradient steps as well as the calibration dataset size. We further improve convergence by combining Teacher and Student activations with a joint distillation loss.
We experimentally show the robustness of our approach by compressing several state-of-the-art small and large Transformers, Mixture-of-Experts, and Mamba LLMs.
Figure~\ref{fig:overview} gives an intuitive overview of the proposed method, which is detailed in Section~\ref{sec:algo} and validated in Section~\ref{sec:xp}.

Our approach addresses some limitations of the compression approaches
presented in Section \ref{sec:llm_compression}. Compared to unstructured pruning
and semi-structured pruning, which result in models with sparse weights and do not achieve speedup without specially designed hardware-specific kernels, our approach, which shrinks
the LLM weights, saves memory across all hardware and leads to speedup, as fewer computations are performed. Structured pruning approaches can be costly: for example, the pruning step of Sheared-LLaMA \citep{xia_sheared_2023} is five times as expensive as standard LM training, according to the authors. In contrast, thanks to local gradient updates, our approach is computationally efficient, allowing us to compress a 47B model within minutes on a single A100 GPU.

There are also some low-rank decomposition approaches that compress models by approximating the pretrained matrices with lower-dimensional ones. Building on the observation that neural networks have low intrinsic dimensions~\citep{li2018measuring}, subsequent work shows that Transformer language models also require lower intrinsic dimensions depending on the language task~\citep{aghajanyan2020intrinsic}. However, other studies show that the weights in some pretrained Transformer models are full-rank, unlike their activations, and propose activation-aware compression methods~\citep{yu_compressing_2023,chen_drone}. Our work is related to these approaches, as we also decompose the matrices of the models. However, we apply a lightweight layer-wise feature distillation objective to better recover from compression and propose different distillation strategies to accelerate convergence and improve performance.

Standard knowledge distillation is data-intensive when the student model's parameters are randomly initialized. Hence, this process is as costly as a standard pretraining stage. Layer-wise distillation can be more suitable in some cases, as intermediate layers may contain certain information, such as speaker identity or phonemes in pretrained speech models \citep{baevski2020wav2vec20frameworkselfsupervised}, \citep{pasad2022layerwiseanalysisselfsupervisedspeech}, \citep{chang2022distilhubertspeechrepresentationlearning}, \citep{chang2021explorationselfsupervisedpretrainedrepresentations}. We leverage the observation that layer activations are low-rank and approximate them with fewer parameters by distilling the Transformer layer activations into low-rank weights
initialized from the SVD of the base weights. This layer-wise distillation approach
enables us to perform local gradient updates, hence making the distillation process cheaper.

\section{Motivation: Low-Rank Activations}
\label{ssec:srank}
A motivation for compressing model weights with SVD would be a low-rank structure in those weights. However, previous studies \citep{chen_drone}, \citep{yu_compressing_2023} have shown that the weights of pretrained Transformer models are often full-rank, while their activations have lower rank.
We empirically confirm this observation on three recent models: Falcon-Mamba, Mistral-v0.1 7B, and Phi-2 3B. There are different measures of rank. In this study, we use a simpler one: stable rank ($srank$) as a proxy measure for rank. The stable rank of a matrix is defined as the following ratio:
\begin{equation}
\text{srank}(X) = \frac{\sum_{i=1}^{r} \sigma_i^2}{\sigma_1^2}
\end{equation}
with $X \in \mathbb{R}^{d_{1} \times d_{2}}$ being the input matrix (activations or weights), $r = \min(d_{1}, d_{2})$, and $\sigma_1 \geq \dots \geq \sigma_r$ the singular values of $X$.

Figure~\ref{fig:srank} shows the stable rank of various weight matrices and layer activations for three models: Falcon-Mamba, Mistral-v0.1 7B, and Phi-2 3B. For all models, activations are significantly lower-rank than weight matrices, which confirms previously published results~\citep{yu_compressing_2023}.

This not only motivates the use of activations for compression, but also explains why applying SVD to weights alone does not work well, as we will show in the experiments. How can we explain the fact that weights are full-rank while activations are low-rank? This is an open question. We know that modern LLMs are pre-trained on large and diverse datasets from the web that contain a great deal of knowledge. Hence, the weights of the model need to be full-rank to capture all this knowledge. However, for a given input from a specific task, only a subset of this knowledge is needed, which may explain the low-rank structure of activations.

\begin{figure*}[h!]
    \centering
    \includegraphics[width=\textwidth]{papers/lillama/stablerank_40p.pdf}
    \caption{\textbf{Activations are low-rank.} Comparison of the stable rank of weights (\texttt{*\_proj}) and layer activations (\texttt{layer\_activations}). Each bar is the average stable rank across all layers.}
    \label{fig:srank}
    \begin{insightbox}{Observation 1: Transformer layer activations are extremely low-rank}
    Transformer layer activations have extremely low stable rank compared to the weights. This extends the previous observations discussed above, which focused on linear activations, to the outputs of entire Transformer layers. Interestingly, the key projection weights also appear to have low stable rank.
    \end{insightbox}
\end{figure*}

\section{Background: Low-Rank Approximation of Transformer Models}
Before presenting our approach, we first introduce two methods for low-rank approximation of Transformer models: low-rank approximation from weights and from features.
\subsection{Low-Rank Approximation from Weights}
Given $W \in \mathbb{R}^{d_{1} \times d_{2}}$, a pretrained weight matrix, and an example $x \in \mathbb{R}^{d_{2}}$, the activations $y \in \mathbb{R}^{d_{1}}$ are computed as $y = Wx$. It is possible to reduce the computations in this operation by using a low-rank approximation of $W$: $y \simeq \Delta Wx = ABx$ with $\Delta W = AB$ the low-rank decomposition of $W$, composed of $A \in \mathbb{R}^{d_{1} \times r}$ and $B \in \mathbb{R}^{r \times d_{2}}$. When the rank $r$ is small enough, the number of parameters $r(d_1+d_2)$ in the low-rank equation is smaller than $d_1d_2$ in the full-rank equation. Estimating the low-rank matrix $\Delta W = AB$ can be formulated as a minimization problem to find the low-rank $\Delta W$ that best approximates $W$:

\begin{equation}
\widehat{\Delta W} = \underset{{\Delta W}}{\mathrm{argmin}} \;\; \|W - {\Delta W}\|_{F}
\end{equation}

It is well-known that this minimization problem can be approached using SVD, a low-rank approximation method that offers the optimal rank-$r$ approximation of a given matrix with respect to the Frobenius norm $\|\cdot\|_{F}$.
SVD decomposes a matrix $W$ into three matrices: $W = USV^{T}$, where $S\in \mathbb{R}^{d_{1} \times d_{2}}$ is a diagonal matrix containing the singular values of $W$ sorted in descending order, and $U \in \mathbb{R}^{d_{1} \times d_{1}}$ and $V \in \mathbb{R}^{d_{2} \times d_{2}}$ are orthonormal matrices. The optimal low-rank approximation of $W \in \mathbb{R}^{d_{1} \times d_{2}}$ can be obtained by keeping the first $r$ singular values, with $r < \min(d_{1}, d_{2})$:

\begin{equation}
\widehat{\Delta W} = (U_{:, :r}S_{:r, :r})V_{:r,:} = AB
\label{eq:svd}
\end{equation}

with $A=U_{:,:r}S_{:r,:r}$ and $B=V_{:r,:}$ where the notation $_{:a, :b}$ refers to the slicing operation.

\subsection{Low-Rank Approximation from Features}
\label{ssec:pca}

As discussed previously, many studies \citep{chen_drone}, \citep{yu_compressing_2023} have shown that the activations (i.e., features) of pretrained transformers have lower rank than the weights. Recent work shows that Transformer activations can be sparsified up to 60\%~\citep{TEAL}.
This motivates using the activations of the linear layers to estimate the low-rank matrices.

Let $\mathcal{D} = \{x_i \in \mathbb{R}^{d_{2}}\}_{1\leq i\leq N}$ be a calibration dataset, and $W \in \mathbb{R}^{d_{1} \times d_{2}}$ the parameters of a linear layer. Finding the low-rank matrix $\Delta W$ that best reproduces the activations $\{y=Wx\}~~~ \forall x\in \mathcal{D}$ involves solving:
\begin{equation}
\widehat{\Delta W} = \underset{{\Delta W}}{\mathrm{argmin}} \;\; \frac{1}{N}\sum\limits_{x \in \mathcal{D}}\|Wx - {\Delta Wx}\|_{F}
\end{equation}
An analytic solution to this minimization problem is the eigendecomposition of the covariance matrix of the activations. We first collect all activations $\mathcal{Y}=\{y=Wx\}_{\forall x\in \mathcal{D}}$. The covariance matrix of the activations can be estimated as $\Sigma = \underset{y \in \mathcal{Y}}{\mathbb{E}}\left[yy^T\right] - \mathbb{E}[y]\mathbb{E}[y]^T$ with $\Sigma \in \mathbb{R}^{d_{1} \times d_{1}}$. Since $\Sigma$ is a diagonalizable matrix, we can apply its eigendecomposition: $\Sigma = USU^{T}$, where $U \in \mathbb{R}^{d_{1} \times d_1}$ contains the eigenvectors of $\Sigma$ and $S \in \mathbb{R}^{d_{1} \times d_{1}}$ is the diagonal matrix containing the eigenvalues sorted in decreasing order. As for the SVD equation above, we can keep the eigenvectors corresponding to the $r$ largest eigenvalues, which gives us $A = U_{:, :r}$ and $B = U_{:, :r}^{T}W$.

Compared to the low-rank matrices in the SVD equation above, the weights $A$ and $B$ here are learned to reproduce the activations of the base weight $W$, thanks to the minimization objective above. In the next section, we highlight the limitations of this approach and introduce our proposed method.

\section{Proposed Approach}
\begin{figure}[t]
    \centering
    \subfloat[Teacher\label{fig:distillation_teacher}]{%
        \includegraphics[width=0.30\linewidth]{papers/lillama/Teacher.pdf}}\hfill
    \subfloat[Student\label{fig:distillation_student}]{%
        \includegraphics[width=0.30\linewidth]{papers/lillama/Student.pdf}}\hfill
    \subfloat[Teacher+Student\label{fig:distillation_joint}]{%
        \includegraphics[width=0.30\linewidth]{papers/lillama/Teacher+Student.pdf}}
    \caption[Accelerating convergence through a joint loss]{\textbf{Accelerating convergence by learning from \protect\colorbox{insightframe}{Teacher} and \protect\colorbox{insightback}{Student} activations through a joint loss.} We compare three distillation strategies: \textbf{(a) Teacher}: the input to the compressed student layer comes from the output of the previous teacher layer; \textbf{(b) Student}: the input to the compressed student layer comes from the output of the previous student layer; \textbf{(c) Teacher+Student}: the compressed student layer is applied separately to the output of the previous teacher layer and to the output of the previous student layer, and the joint loss is the sum of the two corresponding distillation losses.}
    \label{fig:distillaton}
\end{figure}
We next identify and address three potential limitations of the approach presented in Section \ref{ssec:pca}.

\subsection{Non-linear Feature Approximation.}
The analytical solution above is limited to activations from linear layers.
To generalize to non-linear modules, we first observe that this objective can be seen as a feature distillation objective that may be optimized numerically by gradient descent rather than analytically by eigendecomposition. Given the input batch $X \in \mathbb{R}^{d \times b}$ of $b$ examples, we denote the output activations of the $i^{th}$ Teacher module $\mathcal{T}^{(i)}$ as
\begin{equation}
Y^{(i)}= \mathcal{T}^{(i)}(X; \Theta^{(i)})
\end{equation}
where $\Theta^{(i)}$ are the original pretrained matrices of the $i^{th}$ Teacher, and $Y^{(i)} \in \mathbb{R}^{d \times b}$ are its output activations. Similarly, we denote the output activations of the $i^{th}$ Student module $\mathcal{S}^{(i)}$ as:
\begin{equation}
\widehat{Y}^{(i)}= \mathcal{S}^{(i)}(X; \Delta \Theta^{(i)})
\end{equation}
where the Student module is parameterized with the low-rank matrices $\Delta \Theta^{(i)}$ and $\widehat{Y}^{(i)} \in \mathbb{R}^{d \times b}$ are the output activations. We can express the distillation objective of the $i^{th}$ Student module as:
\begin{equation}
\widehat{\Delta  \Theta^{(i)}} = \underset{{\Delta  \Theta^{(i)}}}{\mathrm{argmin}} \;\; \mathcal{L}^{(i)}(Y^{(i)}, \; \widehat{Y}^{(i)})
\end{equation}
where $\widehat{\Delta  \Theta^{(i)}}$ are the estimated low-rank matrices of the $i^{th}$ Student module and $\mathcal{L}^{(i)}(Y^{(i)}, \; \widehat{Y}^{(i)})$ is the loss measuring the distance between the activations of the Teacher and the Student.
We use the combined loss from prior work~\citep{chang2022distilhubertspeechrepresentationlearning} because we observed instability with the $\ell_1$ loss alone:
\begin{align}
\mathcal{L}^{(i)} &= \mathcal{L}^{(i)}_{\ell_1} + \mathcal{L}^{(i)}_{\cos} \\
&= \sum_{t=1}^{b} \left[ \frac{1}{D} \left\| Y^{(i)}_{t} - \widehat{Y}^{(i)}_{t} \right\|_1
- \log \sigma \left( \cos \left( Y^{(i)}_{t}, \widehat{Y}^{(i)}_{t} \right) \right) \right] \notag
\end{align}
where $D$ is the hidden-vector dimension, $\sigma$ is the sigmoid activation and cos$(\cdot, \cdot)$ is the cosine similarity. We propose to approximate the low-rank weights of the Students through gradient descent. In the linear case, this should converge towards the eigendecomposition solution described in Section \ref{ssec:pca}.
However, as we show later, we extend the distillation process to non-linear modules, which requires numerical optimization. We also show that initializing the Student's low-rank parameters with SVD, rather than randomly, improves convergence.

\subsection{Beyond Teacher-only Activations.}
The second potential limitation of the approach described in Section \ref{ssec:pca} and equations above is that the Student and the Teacher modules take the same input $X$, which is the output of the previous layer of the Teacher. This means that the Student is trained from the activations of the Teacher module only, which are not available at inference time. This approach has been used to approximate the linear layers of Transformer models under the name Atomic Feature Mimicking~\citep{yu_compressing_2023}. Let $\widehat{Y}^{(i)}_{T} = \mathcal{S}^{(i)}(Y^{(i-1)}_{T}; \Delta \Theta^{(i)})$ be the output activations of the Student module when fed $Y^{(i-1)}_{T}$, the output activations of the Teacher at the previous layer $i-1$. The loss can then be written as:
\begin{equation}
\mathcal{L}^{(i)}_{T} = \mathcal{L}^{(i)}(Y^{(i)}_{T}, \; \widehat{Y}^{(i)}_{T})
\end{equation}
where $Y^{(i)}_{T}$ are the gold activations of the current Teacher. This approach can lead to fast convergence as the Student modules learn from gold and high-quality activations produced by the Teacher modules. However, it can cause a performance drop during inference as the Teacher activations are not available. An alternative approach is to have the Student module take as input the activations of previous Student modules rather than those of Teacher modules:
\begin{equation}
\mathcal{L}^{(i)}_{S} = \mathcal{L}^{(i)}(Y^{(i)}_{T}, \; \widehat{Y}^{(i)}_{S})
\end{equation}
where $\widehat{Y}^{(i)}_{S} = \mathcal{S}^{(i)}(\widehat{Y}^{(i-1)}_{S}; \Delta \Theta^{(i)})$ denotes the output activations of the current Student module when taking as input the output activations $\widehat{Y}^{(i-1)}_{S}$ of the Student module at the previous layer, and $Y^{(i)}_{T}$ are the gold activations of the corresponding Teacher module.
However, although we found that this loss generally performs
better than the previous one, it leads to slow convergence due to
the errors in the activations of the Student modules. To address both issues,
we propose to combine the two losses:
\begin{equation}
\mathcal{L}^{(i)}_{T+S} = \mathcal{L}^{(i)}_{T} + \mathcal{L}^{(i)}_{S}
\end{equation}

This distillation approach is illustrated in Figure~\ref{fig:distillation_joint}. We don't introduce any hyperparameter for controlling the importance of one loss over another in the joint loss. We leave this for future work.

\subsection{Module-wise Distillation.}
Finally, we observe that the modules need not be limited to linear layers: any sub-part of the model that outputs low-rank activations can be compressed with our method. In our experiments, we distill at the Transformer layer level and leave comparisons with other distillation levels for future work. We also experiment with models other than pretrained Transformer-based language models.

\section{Fast and Memory-Efficient Compression}
\label{sec:algo}

Computing the optimal compression rank of each matrix individually may be costly: while search algorithms can be a solution for smaller models, they can be difficult to scale to LLMs with billions of parameters.
Assuming a target compression rate $N$, we consider three simple strategies: \textbf{uniform}, \textbf{top}, and \textbf{bottom}-first compression. Figure~\ref{fig:strategies} summarizes the strengths and weaknesses of each strategy.

\begin{figure}[h!]
    \centering
    \includegraphics[width=\linewidth]{papers/lillama/strategies.pdf}
    \caption{Illustration of 20\% compression when using the three strategies: bottom, top, and uniform. Green boxes are compressed layers with low-rank matrices while gray boxes are non-compressed ones. $l_{1}$, $l_{2}$, ..., $l_{n-1}$, $l_{n}$ are the layer indices.}
    \label{fig:strategies}
\end{figure}

\paragraph{Uniform Weights Compression}
This approach removes $N$\% of the parameters from each weight matrix in the model. It makes no assumptions about the importance of different layers but is costly because it requires computing the low-rank approximation of all weight matrices. It also requires training all the model's layers, increasing GPU memory and computation requirements through forward and backward passes across the entire model.

\paragraph{Top Layers First Compression}
This approach removes $N$\% of the parameters from the model by prioritizing the top layers. This is more efficient than the uniform approach as it requires computing the low-rank approximation of fewer weight matrices and training fewer layers. However, it still requires loading the full model in GPU memory and performing a forward pass through all the layers, although only a subset is compressed and trained.

\paragraph{Bottom Layers First Compression}
This approach removes $N$\% of the parameters from the model by prioritizing the bottom layers. Compared to the previous two approaches, this one is more efficient as it requires computing the low-rank approximation of fewer weight matrices and training fewer layers. It also allows us to load only a subset of the model into GPU memory and run forward passes through only the first compressed layers, which is not the case for the previous two approaches. As we will show, this approach can compress large models that could not fit into GPU memory. Algorithm~\ref{alg:two} describes the bottom-first compression approach.

\begin{algorithm}[!ht]
\small
\caption{Bottom Layers First Compression Algorithm}\label{alg:two}
\KwIn{$\mathcal{M}$ is the base model; $S$ is the target size of the compressed model; $k$ is the minimum possible rank to set; $m$ is the increment when generating ranks}
\KwResult{Low-rank compressed $\mathcal{M}$}
$\mathcal{R} \gets \emptyset$\;
\For{each weight matrix $W \in \mathbb{R}^{d_{1} \times d_{2}}$ in $\mathcal{M}$}{
    $b \gets \min(d_{1}, d_{2})$\;
    \For{$r = k;\ r \leq b;\ r = r + m$}{
        \If{$r \times (d_{1} + d_{2}) < d_{1} \times d_{2}$}{
            $\mathcal{R} \gets \mathcal{R} \cup \{(r, W)\}$\;
        }
    }
}
sort $\mathcal{R}$ by layer index in increasing order, then by rank in decreasing order\;
\While{$|\mathcal{M}| > S$}{
    get the next $(r, W)$ from $\mathcal{R}$\;
    compute $A, B$ from $W$ using Eq.~\ref{eq:svd} with rank $r$\;
    replace $W$ by $A, B$ in $\mathcal{M}$\;
}
\Return{$\mathcal{M}$}\;
\end{algorithm}

\subsection{Ease of Implementation}
The approach is straightforward to implement. Figure~\ref{fig:pseudocode_impl} shows a simple pseudocode implementation that can be easily integrated into most recent PyTorch models.

\begin{figure}[h!]
    \centering
    \includegraphics[width=\linewidth]{papers/lillama/pseudo-code.pdf}
    \caption{Example of a PyTorch implementation of our approach using the Teacher-only distillation strategy.}
    \label{fig:pseudocode_impl}
\end{figure}

\section{Experiments on Transformers}
\label{sec:xp}
\subsection{Setup}
\paragraph{Models.}
We evaluate our method on several Transformer-based LLMs: a 47B Mixture-of-Experts model (Mixtral-8x7B-v0.1), medium-sized LLMs (Mistral-v0.1 7B and Phi-3 14B), and a smaller model (Phi-2 3B).

\paragraph{Data.}
For calibration, we use 13 million tokens randomly sampled from Slim-Orca \citep{SlimOrca}, an open-source replication of Orca \citep{mukherjee2023orca}. In preliminary experiments, we also tested other datasets such as RedPajama \citep{together2023redpajama}, but observed better results with instruction data.

\paragraph{Implementation Details.}
We implement our approach using version 4.44.2 of the Hugging Face Transformers library~\citep{wolf2020huggingfaces} and PyTorch version 2.4.0+cu121. For all experiments, we use the AdamW optimizer with a learning rate of $8.6e^{-4}$. Note that since our approach is local, each student layer has its own optimizer. This also allows us to perform local gradient updates and avoid storing the entire large PyTorch computation graph in memory. All experiments were conducted on a single A100 GPU, regardless of the model.
When compressing the models using Algorithm~\ref{alg:two}, we set a minimum rank of $k=1024$ for all models except Phi-3 14B, for which we found higher ranks work well, so we set $k=1536$. For all models, we set $m=256$ for the increment to generate ranks.
We apply the Teacher+Student loss, as illustrated in Figure~\ref{fig:distillation_joint}, to all models, as we found that this loss converges faster and generally performs better. Additionally, we use the bottom-first compression strategy, which proved to be effective and less memory-intensive than other strategies.

\paragraph{Evaluation.}
To measure inference speed, we evaluate \textit{prefill throughput} by timing a batch of prompts passed through the entire model and reporting the results in tokens per second. Using lm-evaluation-harness \citep{eval-harness}, we evaluate zero-shot performance on 9 downstream tasks: TruthfulQA (TruQA) \citep{lin-etal-2022-truthfulqa}, Social IQA (SIQA) \citep{sap2019social}, LogiQA \citep{liu2020logiqa}, WinoGrande (WinoG) \citep{sakaguchi2019winogrande}, ARC Easy (ARC-E) \citep{arc}, ARC Challenge (ARC-C) \citep{arc}, BoolQ \citep{boolq}, PIQA \citep{piqa}, and OpenBookQA (OBQA) \citep{obqa}. We compare non-compressed (0\%) and compressed variants.
We provide in Table~\ref{tab:task_metric_map} the metrics we used to evaluate the models in Section \ref{sec:xp}.
\begin{table}[!ht]
\centering
    \begin{tabular}{ll}
    \toprule
    \textbf{Task} & \textbf{Metric} \\ 
    \midrule
    arc\_challenge   & acc\_norm \\ 
    arc\_easy        & acc\_norm \\ 
    piqa             & acc\_norm \\ 
    social\_iqa      & acc \\ 
    logiqa           & acc\_norm \\ 
    truthfulqa\_mc2  & acc \\ 
    winogrande       & acc \\ 
    boolq            & acc \\ 
    openbookqa       & acc\_norm \\ 
    \bottomrule
    \end{tabular}
\caption{Benchmarks and corresponding metrics that were used to evaluate the models.}
\label{tab:task_metric_map}
\end{table}

\subsection{Results}
\paragraph{Zero-shot evaluation under different compression ratios.}
Table~\ref{tab:main_result} reports zero-shot results for base models and 20\%, 25\%, and 30\% compressed models, showing that the method remains robust as compression increases.

\paragraph{Low-rank models are lighter and faster.}
Table~\ref{tab:mem-speed} shows memory savings (VRAM) and speedup of compressed models. We can see that 20\% compressed models are lighter and up to 20\% faster.

\paragraph{Our approach can efficiently create small language models.}
We compressed Phi-2 3B to 1.7B (40\% parameter reduction) and compare its performance to recent models of similar size in Table~\ref{tab:smollm}.

\paragraph{Compressed models recover well with fine-tuning.}
Table~\ref{tab:ft_result} shows that fine-tuning these compressed models leads to performance recovery.

\begin{insightbox}{Observation 2: Low-rank compressed LLMs can continue to improve}
As shown by the Mistral fine-tuning results in Table~\ref{tab:ft_result}, the bottleneck representation created by the low-rank structure does not seem to prevent the model from improving. With longer fine-tuning, it could potentially recover most of the lost performance.
\end{insightbox}

\begin{table*}[t]
    \resizebox{\textwidth}{!}{%
    \setlength{\tabcolsep}{3pt}
    \begin{tabular}{cccccccccccc}
        \toprule
        \textbf{Model} & \textbf{Reduction} & \textbf{TruQA} & \textbf{SIQA} & \textbf{LogiQA} & \textbf{WinoG} & \textbf{ARC-E} & \textbf{ARC-C} & \textbf{BoolQ} & \textbf{PIQA} & \textbf{OBQA} & \textbf{Average} \\
        \midrule 
        \multirow{4}*{\centering Phi-3 14B}
        & 0\% & 57.63 & 57.06 & 37.94 & 76.01 & 81.36 & 61.60 & 88.56 & 81.34 & 50.40 & \textbf{65.77} \\
        & 20\% & 57.88 & 50.26 & 34.10 & 72.53 & 83.54 & 59.22 & 85.32 & 80.30 & 49.20 & \textbf{63.59} \\
        & 25\% & 56.37 & 51.38 & 32.26 & 70.56 & 80.98 & 57.51 & 83.88 & 79.76 & 47.00 & \textbf{62.19} \\
        & 30\% & 55.32 & 50.41 & 31.18 & 70.32 & 78.70 & 55.55 & 84.10 & 76.99 & 45.00 & \textbf{60.98} \\
        \midrule
        \multirow{4}*{\centering Phi-2 3B}
        & 0\% & 44.40 & 55.42 & 30.57 & 76.16 & 78.16 & 54.18 & 83.21 & 79.11 & 51.20 & \textbf{61.38} \\
        & 20\% & 46.04 & 52.56 & 29.80 & 71.51 & 72.81 & 46.50 & 75.05 & 76.55 & 45.60 & \textbf{57.38} \\
        & 25\% & 44.55 & 51.64 & 30.88 & 71.35 & 71.93 & 45.99 & 79.79 & 74.43 & 43.20 & \textbf{57.08} \\
        & 30\% & 44.12 & 51.59 & 30.41 & 70.72 & 69.99 & 44.37 & 78.17 & 74.54 & 43.20 & \textbf{56.35} \\
        \midrule
        \multirow{4}*{\centering Mistral-v0.1 7B}
        & 0\% & 42.60 & 46.57 & 29.80 & 73.80 & 79.55 & 53.92 & 83.70 & 82.10 & 44.00 & \textbf{59.56} \\
        & 20\% & 44.14 & 45.50 & 28.26 & 66.69 & 73.11 & 46.33 & 77.00 & 77.37 & 41.20 & \textbf{55.51} \\
        & 25\% & 40.52 & 44.93 & 30.11 & 68.11 & 65.07 & 40.61 & 78.44 & 73.99 & 40.20 & \textbf{53.55} \\
        & 30\% & 40.08 & 44.37 & 30.57 & 66.93 & 65.53 & 40.61 & 77.98 & 73.18 & 39.00 & \textbf{53.14} \\
        \bottomrule
    \end{tabular}}
    \caption{\textbf{Zero-shot evaluation results} for the base non-compressed models (0\%), and the 20\%, 25\%, and 30\% compressed models.}
    \label{tab:main_result}
\end{table*}

\begin{table*}
    \resizebox{1.06\linewidth}{!}{%
    \begin{tabular}{cccc|ccccccc}
    \toprule
       \textbf{Model} & \textbf{Reduction}  & \textbf{Model Size} & \textbf{VRAM} &  s = 512 & s = 1024 &  s = 2048 & s = 4096 & s = 8192 & s = 16384 \\
    \midrule
    \multirow{2}*{\centering Phi-3 14B} 
       & 0\%      & 14B &  28 GB & 7171 t/s & 7216 t/s & 7197 t/s & 7057 t/s & 7010 t/s & 7036 t/s \\
       & 20\%     & 11B &  22.8 GB & 8622 t/s & 8686 t/s & 8509 t/s & 8349 t/s & 8268 t/s & 8269 t/s \\
    \midrule
       \multirow{2}*{\centering Mixtral-8x7B-v0.1 47B} 
       & 0\%      & 47B & OOM & OOM & OOM & OOM & OOM & OOM & OOM \\
       & 20\%     & 37B &  73.8 GB & 6515 t/s & 8143 t/s & 8444 t/s & OOM & OOM & OOM \\
    \midrule
       \multirow{2}*{\centering Phi-2 3B} 
       & 0\%      & 2.8B  &  6.8 GB & 29949 t/s & 30895 t/s & 30184 t/s & 30403 t/s & 26636 t/s & 21499 t/s \\
       & 20\%     & 2.2B  &  5.7 GB & 32451 t/s & 34561 t/s & 34276 t/s & 32479 t/s & 30351 t/s & 23735 t/s \\
    \midrule
       \multirow{2}*{\centering Mistral-v0.1 7B} 
       & 0\%      & 7.2B  &  15.2 GB & 13385 t/s & 13537 t/s & 13650 t/s & 13265 t/s & 12603 t/s & 12399 t/s \\
       & 20\%     & 5.8B  &  12.6 GB & 15416 t/s & 15832 t/s & 15947 t/s & 15446 t/s & 14589 t/s & 14304 t/s \\
    \bottomrule
    \end{tabular}}
    \caption{\textbf{Compressed models save memory and speed up computation.} Inference speed was measured on a single A100 GPU, using the whole test split of Wikitext2 with a batch size of 4 and varying the sequence length from 512 to 16384 tokens. All the models were loaded in bf16 and used Flash-Attention2. OOM stands for out of memory.}
    \label{tab:mem-speed}
\end{table*}

\begin{table*}[t]
    \resizebox{1.06\textwidth}{!}{%
    \setlength{\tabcolsep}{3pt}
    \begin{tabular}{lcccccccccccc}
        \toprule
        \textbf{Model} & \textbf{TruQA} & \textbf{SIQA} & \textbf{LogiQA} & \textbf{WinoG} & \textbf{ARC-E} & \textbf{ARC-C} & \textbf{BoolQ} & \textbf{PIQA} & \textbf{OBQA} & \textbf{Average} \\
        \midrule 
        StableLM-2 1.6B & 38.90 & 48.52 & 26.73 & 63.30 & 68.39 & 38.57 & 74.80 & 76.93 & 39.00 & \textbf{52.79} \\
        Qwen-1 1.8B & 38.09 &  45.19 & 31.80 & 59.12 & 58.33 & 34.98 & 65.93 & 73.23 & 33.60 & \textbf{48.92} \\
        Qwen-2 1.5B &  45.93 & 45.85 & 31.18 & 66.22 & 60.56 & 36.09 & 72.26 & 75.35 & 36.40 & \textbf{52.20} \\
        \textbf{Lillama} (Phi-2 1.7B) & 44.08 & 47.80 & 25.65 & 68.27 & 63.22 & 38.82 & 76.02 & 72.36 & 39.20 & \textbf{52.82} \\
        \bottomrule
    \end{tabular}}
    \caption{\textbf{40\% compressed Phi-2 3B competes with models of similar size while being compressed using only 13M tokens.} We compress Phi-2 3B to 1.7B (by removing 40\% of the parameters) and compare its zero-shot performance to that of other recent language models of similar sizes.}
    \label{tab:smollm}
\end{table*}

\begin{table*}[t]
    \resizebox{1.06\textwidth}{!}{%
    \setlength{\tabcolsep}{3pt}
    \begin{tabular}{cccccccccccc}
        \toprule
        \textbf{Model} & \textbf{Reduction} & \textbf{TruQA} & \textbf{SIQA} & \textbf{LogiQA} & \textbf{WinoG} & \textbf{ARC-E} & \textbf{ARC-C} & \textbf{BoolQ} & \textbf{PIQA} & \textbf{OBQA} & \textbf{Average} \\
        \midrule 
        \multirow{3}*{\centering Phi-2 3B} 
        & 0\% & 44.40 & 55.42 & 30.57 & 76.16 & 78.16 & 54.18 & 83.21 & 79.11 & 51.20 & \textbf{61.38} \\
        & 40\% & 44.08 & 47.80 & 25.65 & 68.27 & 63.22 & 38.82 & 76.02 & 72.36 & 39.20 & \textbf{52.82} \\
        & 40\% {\scriptsize FT} & 45.43 & 49.49 & 31.80 & 70.09 & 60.86 & 37.29 & 67.68 & 73.23 & 42.00 & \textbf{53.10} \\
        \midrule
        \multirow{3}*{\centering Mistral-v0.1 7B} 
        & 0\% & 42.60 & 46.57 & 29.80 & 73.80 & 79.55 & 53.92 & 83.70 & 82.10 & 44.00 & \textbf{59.56} \\
        & 40\% & 41.69 & 43.65 & 30.57 & 62.35 & 61.66 & 34.56 & 75.23 & 71.49 & 34.80 & \textbf{50.63} \\
        & 40\% {\scriptsize FT} & 41.60 & 52.25 & 29.65 & 66.61 & 66.75 & 39.93 & 79.39 & 74.37 & 38.29 & \textbf{54.32}\\
        \bottomrule
    \end{tabular}}
    \caption{\textbf{Compressed models recover well when fine-tuned.} Zero-shot performance of the base non-compressed models (0\%) and the 40\% compressed models, without and with fine-tuning (FT). Models were fine-tuned on the whole Slim-Orca dataset (191 million tokens).}
    \label{tab:ft_result}
\end{table*}


\section{Other Architectures and Modalities}
In this section, we show the generalizability of our approach by evaluating it on the Mamba architecture for text \citep{gu2024mambalineartimesequencemodeling} and on the speech modality with Whisper \citep{radford2022robustspeechrecognitionlargescale}.
We further evaluate it on a lower-resource African language later in this chapter.
\subsection{Mamba Architecture}

The attention mechanism in the Transformer is limited by its increasing complexity as a function of the input sequence length.
Therefore, Linear Attention Models have been proposed as alternatives to Transformers, in particular State Space Models, such as the Mamba architecture \citep{gu2024mambalineartimesequencemodeling}, which have shown promising results.

\paragraph{Method.} We tested our approach on two Mamba architectures: Mamba
3B \citep{gu2024mambalineartimesequencemodeling} (\url{https://huggingface.co/state-spaces/mamba-2.8b-hf}) and Falcon-Mamba 7B (\url{https://huggingface.co/tiiuae/falcon-mamba-7b}).
We compressed these models by 20\% using the same dataset and hyperparameters as in previous experiments with Transformers, with a minimum rank set to 1024 for both models.

\paragraph{Results.} Table~\ref{tab:ssm} shows that, as for Transformers, Mamba-based language models can be efficiently compressed. Interestingly, the 20\% compressed Mamba 3B maintains 99\% of the performance.
Table~\ref{tab:mem-speed-mamba-ch2} further shows that compressed Mamba models are lighter and faster than their original versions.
\begin{table*}[h!]
    \centering
    \resizebox{1.02\textwidth}{!}{%
    \setlength{\tabcolsep}{3pt}
    \begin{tabular}{lccccccccccc}
        \toprule
        \textbf{Model} & \textbf{Reduction} & \textbf{TruQA} & \textbf{SIQA} & \textbf{LogiQA} & \textbf{WinoG} & \textbf{ARC-E} & \textbf{ARC-C} & \textbf{BoolQ} & \textbf{PIQA} & \textbf{OBQA} & \textbf{Average} \\
        \midrule
        \multirow{2}*{\centering Falcon-Mamba 7B} 
        & 0\% & 53.40 & 52.66 & 31.34 & 74.66 & 81.73 & 58.96 & 83.12 & 81.56 & 48.60 & \textbf{62.89} \\
        & 20\% & 52.34 & 50.10 & 30.72 & 66.69 & 76.73 & 49.23 & 77.22 & 78.02 & 44.80 & \textbf{58.43} \\
        \midrule
        \multirow{2}*{\centering Mamba 3B} 
        & 0\% & 35.88 & 43.24 & 26.88 & 63.38 & 64.02 & 36.26 & 65.63 & 75.90 & 39.40 & \textbf{50.07} \\
        & 20\% & 37.50 & 42.43 & 27.04 & 61.40 & 62.16 & 35.58 & 64.89 & 73.29 & 40.00 & \textbf{49.37} \\
        \bottomrule
    \end{tabular}
    }
    \caption{\textbf{Mamba-based language models can be efficiently compressed}. Zero-shot evaluation results for the non-compressed (0\%) and for the compressed Mamba-based language models (20\%).}
    \label{tab:ssm}
\end{table*}

\begin{table*}
    \centering
    \resizebox{\linewidth}{!}{%
    \setlength{\tabcolsep}{3pt}
    \begin{tabular}{ccccc|ccccc}
    \toprule
       \textbf{Model} & \textbf{Reduction} & \textbf{Model Size} & \textbf{VRAM} & s = 512 & s = 1024 & s = 2048 & s = 4096 & s = 8192 & s = 16384 \\
    \midrule
    \multirow{2}*{\centering Falcon-Mamba 7B}
       & 0\%      & 7.27B & 15.3 GB & 8725 t/s & 9346 t/s & 9390 t/s & 9695 t/s & 9908 t/s & 9486 t/s \\
       & 20\%     & 5.82B & 12.6 GB & 10227 t/s & 10427 t/s & 10477 t/s & 10814 t/s & 11108 t/s & 10647 t/s \\
    \midrule
    \multirow{2}*{\centering Mamba 3B}
       & 0\%      & 2.77B & 6.7 GB & 17994 t/s & 18556 t/s & 18819 t/s & 19618 t/s & 19231 t/s & 18602 t/s \\
       & 20\%     & 2.22B & 5.6 GB & 19013 t/s & 19993 t/s & 20378 t/s & 21218 t/s & 20808 t/s & 20100 t/s \\
    \bottomrule
    \end{tabular}}
    \caption{\textbf{Compressing Mamba models saves memory and speeds up computation.} Inference speed was measured on a single A100 GPU, using the whole test split of Wikitext2 with a batch size of 4 and by varying the sequence length from 512 to 16384 tokens. All models were loaded in bf16 and used Flash-Attention2.}
    \label{tab:mem-speed-mamba-ch2}
\end{table*}


\subsection{Speech modality: Whisper}

We next evaluate our approach on a state-of-the-art speech model: \textit{whisper-medium.en} \citep{radford2022robustspeechrecognitionlargescale}, a 764M-parameter encoder-decoder speech recognition model trained on English.

\paragraph{Method.} We take inspiration from the deep encoder,
shallow decoder approach \citep{kasai2021deepencodershallowdecoder}, where the
authors observe that the decoder in encoder-decoder models does not need to be very deep to achieve good performance.
Therefore, we keep the encoder unchanged and remove 70\% of the decoder's parameters. We then approximate the low-rank matrices using the same method and hyperparameters described in Section \ref{sec:xp}.
\paragraph{Data.} For calibration, we used a random 60\% sample of examples from Librispeech-train-100 \citep{librispeech}. We also fine-tuned the compressed model on a subset of 10 hours of speech randomly sampled from Librispeech-train-100.

\paragraph{Results.} We evaluate the models' Word Error Rate (WER), memory footprint and latency on an out-of-distribution test dataset: the test split of the English FLEURS dataset \citep{fleurs}.

\begin{table}[!ht]
    {\small{
    \begin{tabular}{ccccc}
    \toprule
        \textbf{Reduction} &\textbf{Size} & \textbf{WER ($\downarrow$)}  & \textbf{VRAM} & \textbf{Speed} ($\uparrow$) \\
    \midrule
        0\% & 764M & 25.40 & 3.5GB & 2395t/s\\
        37\%  & 485M & 30.84 & 2.4GB & 2717t/s \\
        0\% {\scriptsize FT} & 764M & 9.20 & 3.5GB & 2395t/s \\
        37\% {\scriptsize FT} & 485M & 12.14 & 2.4GB & 2717t/s \\ 
    \bottomrule
    \end{tabular}
    }}
    \caption{\textbf{The 37\% compressed Whisper is faster and 1.46$\times$ lighter.} The WER is computed on the English test subset of the FLEURS dataset. The speed (tokens/second) corresponds to the transcription of 400 tokens from a dataset containing 2,048 speech examples, using a batch size of 128. A single A100 GPU is used and the model is loaded in FP32, without Flash-Attention. We also give the GPU Memory (VRAM) occupied by each model.}
    \label{tab:wer}
\end{table}

Table~\ref{tab:wer} gives the WER for compressed and uncompressed models. We can see that a 37\% compressed Whisper maintains 82\% of the performance.

\section{Analysis}

\begin{figure}[ht]
    \centering
    \includegraphics[width=\linewidth]{papers/lillama/rank_acc-1.pdf}
    \caption{\textbf{Higher ranks generally give better results.} Mean accuracy over our test benchmarks as a function of the minimum rank $k$ in Algorithm~\ref{alg:two}}
    \label{fig:ranks}
\end{figure}

\begin{figure}[ht]
    \centering
    \includegraphics[width=\linewidth]{papers/lillama/w_wo_svd-2.pdf}
    \caption{\textbf{SVD initialization improves sample efficiency}. Convergence when initializing low-rank weights randomly (without SVD) or with SVD.}
    \label{fig:w_wo_svd}
\end{figure}


In this section, we give some ablation studies. In all experiments, the models are compressed by 20\%.


\subsection{Which Distillation Loss Should Be Used?}
\begin{figure*}[h!]
    \centering
        \includegraphics[width=1.0\textwidth]{papers/lillama/losses_convergence.pdf}
        \caption{\textbf{The joint loss generally converges better}. Convergence of the three losses illustrated in Figure~\ref{fig:distillaton}, evaluated using perplexity on the Wikitext2 test corpus during distillation.}
        \label{fig:losses}
\end{figure*}

\begin{table}[h!]
    \centering
    \begin{tabular}{lccc}
    \toprule
    \textbf{Model} & \textbf{Teacher} & \textbf{Student} & \textbf{Tea+Stu} \\
    \midrule
    Phi-2 3B & 57.38 & 56.44 & 57.38 \\
    Phi-3 14B & 62.54 & 63.20 & 63.36 \\
    Mistral-v0.1 7B & 51.25 & 55.46 & 55.51 \\
    Mixtral-8x7B 47B & 58.87 & 60.49 & 60.19 \\
    \midrule
    \textbf{Average} & 57.51 & 58.90 & \textbf{59.11} \\
    \bottomrule
    \end{tabular}
    \caption{\textbf{The joint loss generally performs better}: Zero-shot performance of the three losses using the same setup for each model, as described in Section~\ref{sec:xp}.}
    \label{tab:losses}
\end{table}

Table~\ref{tab:losses} shows that the joint loss depicted in Figure~\ref{fig:distillation_joint} generally produces the best results. Although the differences between the Student and Teacher+Student losses may not seem significant, Figure~\ref{fig:losses} shows that the Teacher+Student loss leads to faster convergence than the Student loss alone. It also generally results in a lower final perplexity than using the Teacher loss only.

\begin{insightbox}{Observation 3: Combining Teacher and Student activations improves average performance}
Combining Teacher and Student activations in a joint loss achieves the highest average zero-shot score across the evaluated models (Table~\ref{tab:losses}). It also converges faster than the Student-only loss and generally reaches a lower final perplexity than the Teacher-only loss (Figure~\ref{fig:losses}).
\end{insightbox}

As shown in Figure~\ref{fig:losses}, the Teacher+Student loss combines the best of both worlds: fast convergence due to the Teacher activations and high performance at inference thanks to the Student activations. We illustrate this in Table~\ref{tab:logs} for the first 2M training tokens. It shows that the Teacher+Student loss converges faster than the Student-only loss. While the Teacher-only loss also converges fast, it reaches a higher plateau and does not converge beyond that point.

\begin{table}[!ht]
    {\scriptsize{
    \begin{tabular}{lcccc}
    \toprule
    \textbf{Model} & \textit{Seen Tokens} & \textbf{Teacher} & \textbf{Student} & \textbf{Tea+Stu} \\
    \midrule
    \multirow{5}{*}{Phi-2 3B}
                              & 0       & 7185.31  & 7185.31  & 7185.31 \\
                              & 532480  & 20.74    & 36.56    & 20.34 \\
                              & 1056768 & 17.69    & 25.11    & 17.26 \\
                              & 1559879 & 16.00    & 18.97    & 15.92 \\
                              & 2000506 & 15.21    & 18.19    & 15.51 \\
    \midrule
    \multirow{5}{*}{Mistral-v0.1 7B}
                              & 0       & 12365.44 & 12365.44 & 12365.44 \\
                              & 532480  & 43.79    & 270.16   & 114.02 \\
                              & 1056768 & 19.07    & 53.90    & 21.90 \\
                              & 1579121 & 18.56    & 22.40    & 13.64 \\
                              & 2056609 & 18.55    & 15.72    & 13.11 \\
    \midrule
    \multirow{5}{*}{Phi-3 14B}
                              & 0       & 46362.82 & 46362.82 & 46362.82 \\
                              & 532480  & 9.77     & 12.94    & 8.89 \\
                              & 1056768 & 8.33     & 8.76     & 7.99 \\
                              & 1580799 & 7.70     & 7.85     & 7.44 \\
                              & 2069488 & 7.61     & 7.77     & 7.34 \\
    \bottomrule
    \end{tabular}
    }}
    \caption{Comparing the Wikitext perplexity convergence of the three losses for 2M training tokens.}
    \label{tab:logs}
\end{table}

\subsection{How Should the Rank Be Chosen?}
In our bottom-first compression strategy, a small minimum rank in
Algorithm~\ref{alg:two} compresses the bottom layers more aggressively and leaves
more upper layers uncompressed, hence saving GPU memory. However, a
small rank may also destroy knowledge and the model will struggle to recover
during the distillation phase. Figure~\ref{fig:ranks} compares 4 rank
values for Mistral-v0.1 7B and Phi-3 14B and shows that higher ranks
lead to better performance, but are also more costly as more layers are distilled.

\subsection{Is SVD initialization necessary?}
Figure~\ref{fig:w_wo_svd} shows the perplexity curves for Phi-2 3B and Phi-3 14B when the low-rank matrices are
initialized with SVD or randomly. SVD initialization enables faster
convergence, with 8 million tokens being sufficient, i.e., roughly
4k examples of 2048 tokens.

\subsection{Cost.}
All models in Section~\ref{sec:xp} are compressed in less than one hour on a single A100 GPU. The precise cost depends on the number of layers distilled, which is determined by the compression strategy (bottom-first, top-first, or uniform) and the minimum rank $k$ in Algorithm~\ref{alg:two}. Figure~\ref{fig:strategies} intuitively illustrates the three compression strategies evaluated in our method: uniform, bottom-first, and top-first, explaining their respective advantages and drawbacks. We present here empirical results on differences in efficiency and performance using the same hyperparameters as in Section~\ref{sec:xp}.

\paragraph{Differences in running time and memory.}
The bottom-first compression strategy allows compressing all models, including Mixtral-8x7B-v0.1 47B. While the top-first strategy can compress Phi-3 14B, the uniform compression strategy fails to compress Phi-3 14B and Mixtral-8x7B-v0.1 (47B). Table~\ref{tab:ctime-ch2} gives the distillation time for each strategy and each model.

\begin{table}[!ht]
    {\small{
    \begin{tabular}{lccc}
    \toprule
    \textbf{Model} & \textbf{Bottom} & \textbf{Top} & \textbf{Uniform} \\
    \midrule
    Phi-2 3B & 25min & 32min & 50min \\
    Mistral-v0.1 7B & 26min & 46min & 105min \\
    Phi-3 14B & 30min & 50min & OOM \\
    Mixtral-8x7B-v0.1 47B & 47min & OOM & OOM \\
    \bottomrule
    \end{tabular}
    }}
    \caption{Compression times for the three compression strategies at 20\% parameter reduction on a single A100 GPU. For the bottom-first and top-first compression strategies, we used the same hyperparameters as in Section~\ref{sec:xp}.}
    \label{tab:ctime-ch2}
\end{table}

\begin{insightbox}{Observation 4: Bottom-first compression is faster and more memory efficient}
Bottom-first compression reduces computation and GPU memory requirements by processing only the required bottom layers. It is the fastest strategy across all evaluated models (Table~\ref{tab:ctime-ch2}). It is also the only strategy that compresses Mixtral 47B on a single A100 GPU, taking 47 minutes.
\end{insightbox}

\paragraph{Differences in performance.}
Table~\ref{tab:cacc-ch2} presents the mean accuracies for each compression strategy. The bottom-first compression strategy performs well and is also more memory efficient. The top-first strategy does not improve accuracy while having higher time and memory complexity compared to the bottom-first approach.

\begin{table}[!ht]
    {\small{
    \begin{tabular}{lccc}
    \toprule
    \textbf{Model} & \textbf{Bottom} & \textbf{Top} & \textbf{Uniform} \\
    \midrule
    Phi-2 3B & 57.38 & 55.71 & 58.21 \\
    Mistral-v0.1 7B & 55.51 & 55.63 & 54.25 \\
    Phi-3 14B & 63.59 & 60.80 & OOM \\
    Mixtral-8x7B-v0.1 47B & 60.19 & OOM & OOM \\
    \bottomrule
    \end{tabular}
    }}
    \caption{Average zero-shot performance for the three compression strategies at 20\% compression. For the bottom-first and top-first compression strategies, we used the same hyperparameters as in Section~\ref{sec:xp}.}
    \label{tab:cacc-ch2}
\end{table}

\begin{table*}[h!]
    \centering
    \resizebox{\textwidth}{!}{%
    \begin{tabular}{cccccccc}
        \toprule
        \textbf{Reduction} & \textbf{Method} & \textbf{PIQA} & \textbf{WinoG} & \textbf{HellaSwag} & \textbf{ARC-E} & \textbf{ARC-C} & \textbf{Avg.} \\
        \midrule
        \multirow{2}{*}{0\%}
        & SliceGPT & 79.11 & 75.77 & 73.83 & 78.32 & 54.18 & 72.24 \\
        & Reproduced & 79.11 & 76.01 & 73.60 & 78.41 & 54.35 & 72.30 \\
        \midrule
        \multirow{2}{*}{24\%}
        & SliceGPT & 74.05 & 62.12 & 53.31 & 67.26 & 39.42 & 59.23 \\
        & \textbf{Lillama} & \textbf{74.76} & \textbf{69.77} & \textbf{60.60} & \textbf{72.10} & \textbf{47.10} & \textbf{64.86} \\
        \bottomrule
    \end{tabular}
    }
    \caption{\textbf{Comparison of our method with SliceGPT on zero-shot task evaluation for Phi-2 3B.} We give the evaluation scores before compression (0\%), both as reported by SliceGPT and as reproduced by us. Both methods use the Alpaca calibration dataset for compression.}
    \label{tab:comparisons}
\end{table*}

\subsection{Examples of generated texts}
We also evaluated the compressed models by observing the text they generated. In general, at 20\% compression of Phi-2 3B, the generation ability and knowledge remained nearly intact. At 40\% compression, the model still retained its general knowledge and generation abilities but tended to be more repetitive and verbose at times. We recommend users fine-tune the 40\% compressed Phi-2 3B on their specific tasks, such as retrieval-augmented generation. Tables~\ref{tab:sj} and~\ref{tab:cw} give some generated texts from 20\% and 40\% compressed Phi-2 3B.

\begin{table*}
\scriptsize
\centering
\setlength{\tabcolsep}{3pt}
\begin{tabular}{ p{0.31\textwidth} | p{0.31\textwidth} | p{0.31\textwidth} }
\toprule
Base Model & -20\% & -40\% \\
\midrule
    Steve Jobs was a young man who had a passion for technology and design. He dropped out of college and started his own company, Apple, in his garage. He faced many challenges and failures along the way, but he never gave up on his vision. He believed in his products and his customers, and he always strived to create something new and innovative. He was not afraid to take risks and to challenge the status quo. He was a leader who inspired others to follow his example. He was a visionary who changed the world. &  Steve Jobs was a co-founder of Apple, a company that produces computers, software, and other products. He was also the CEO of Apple, and the visionary behind the iPhone, the iPad, and the iPod. He was known for his innovation, creativity, and leadership. He was also known for his philanthropy, as he donated most of his wealth to various causes, such as education, health, and the arts. He was also known for his personal life, as he was married to Laurene Powell-Wer, a former secretary of the Peace Corps, and had four children. He died in 2011 at the age of 56, after suffering from pancreatic cancer.

Questions:

- What was Steve Jobs' occupation?
- What were some of the products that Apple produced?
- What were some of the causes that Steve Jobs donated to?
- How many children did Steve Jobs have?
- When and how did Steve Jobs die?

Answers:

- Steve Jobs' occupation was a co-founder of Apple and the CEO of Apple.
- Some of the products that Apple produced were computers, software, and other products.
- Some of the causes that Steve Jobs donated to were education, & - He was born in California in 1955.

- He dropped out of college and started working at a computer company.

- He became the co-founder of Apple.

- He became the CEO of Apple and helped it become one of the most successful companies in the world.

- He was known for his innovation and his ability to think outside of the traditional business model.

- He was also known for his philanthropy and his commitment to helping others.

- He died in 2011.
\\
\bottomrule
\end{tabular}
\caption{Steve Jobs Story by the 20\% and 40\% compressed Phi-2 3B models, without \textit{any} fine-tuning.
\textbf{Prompt: "Here the story of Steve Jobs:"}}
\label{tab:sj}
\end{table*}

\begin{table*}
\scriptsize
\centering
\setlength{\tabcolsep}{3pt}
\begin{tabular}{ p{0.31\textwidth} | p{0.31\textwidth} | p{0.31\textwidth} }
\toprule
Base Model & -20\% & -40\% \\
\midrule
The Civil War was caused by a number of factors, including economic and social differences between the North and the South, as well as disagreements over states' rights and slavery. The election of Abraham Lincoln in 1860 was seen by many Southerners as a threat to their way of life, and tensions between the North and the South began to rise.
The war began in 1861 when Confederate forces attacked Fort Sumter in South Carolina. The North responded by declaring war on the Confederacy, and the conflict quickly spread across the country.
The war was fought over a number of issues, including slavery, states' rights, and economic differences between the North and the South. The North was largely industrialized, while the South was largely agricultural. This led to tensions over tariffs and other economic issues.
The war was also fought over the issue of slavery. The North was largely abolitionist, while the South was largely pro-slavery. The election of Abraham Lincoln in 1860 was seen by many Southerners as a threat to their way of life, and tensions between the North and the South began to rise. & The American Civil War was a conflict that took place from 1861 to 1865. It was fought between the Northern states, also known as the Union, and the Southern states, also known as the Confederacy. The main cause of the war was the issue of slavery. The Southern states wanted to keep slavery, while the Northern states wanted to abolish it. This led to tensions and eventually led to the war.

What were the major battles of the Civil War? Here is the story.

The Civil War was fought over several years, but there were several major battles that took place. The most famous battles were the Battle of Gettysburg, the Battle of Antieta, and the Battle of Bull Run. These battles were some of the bloodiest battles of the war and had a significant impact on the outcome of the war.

What was the Emancipation Proclamation? Here is the story.

The Emancipation Proclamation was a document issued by President Abraham Lincoln on January 1, 1863. It declared that all slaves in Confederate territory were to be set free. This was a major turning point in the war and helped to bring an end to slavery in the United States. & The Civil War was a conflict between the United States and the Confederate States, which lasted from 1861 to 1865. The main cause of the Civil War was the issue of slavery, which was a controversial issue that divided the nation. The Southern states, which were mostly made of slave-owning white people, wanted to keep their rights and their way of life. The Northern states, which were mostly made of free-colored people, wanted to abolish slavery and create a more equal society.

The Civil War began when the Confederate States declared their secession from the United States, which means they wanted to break away from the country. The United States responded by sending troops to the South, which started the war. The war lasted for four years, and it was one of the most important events in American history. The war ended with the defeat of the Confederate States, and the United States became a unified nation.

The Civil War had many consequences, both positive and negative. Some positive consequences were that the United States became more united and more democratic, and that the issue of slavery was finally resolved. Some negative consequences were that many people died, and that the war caused a lot of damage and destruction.
\\
\bottomrule
\end{tabular}
\caption{The Civil War Story by the 20\% and 40\% compressed models, without \textit{any} fine-tuning. \textbf{Prompt: "What is the cause of the Civil War? Here is the story."}}
\label{tab:cw}
\end{table*}

\subsection{Comparisons.}
Table~\ref{tab:comparisons} shows that our approach outperforms SliceGPT \citep{ashkboos_slicegpt_2024}, a one-shot compression method for LLMs.
SliceGPT reports a ``slicing'' proportion, which does not directly correspond to the percentage reduction in the model's parameter count.
For example, a 30\% slicing proportion for Phi-2 3B corresponds to a 24\% reduction in its parameter count.

\subsection{Can this approach work for low-resource and under-represented languages?}
\label{sec:hausa-ch2}

For most human languages, there is not enough data (generally less than 1B tokens) to pretrain large LLMs. In such cases, data-efficient compression approaches are preferred. We experimented with compressing a tiny language model for Hausa, a low-resource language for which there are not billions of tokens available for continued pretraining. We used InkubaLM~\citep{tonja2024inkubalm}, a 422M-parameter language model designed for five low-resource languages. Since 60\% of the parameters are in the embeddings, we focused on compressing the input embedding $W_{e}$ and the prediction head $W_{o}$. We will discuss later on the use of low-rank structure in the LM head, as it appears to contradict some empirical observations from previous studies. For now, let's focus on compressing the input and output modules by applying our approach through local distillation of the low-rank embeddings and prediction head:

\[
\widehat{\Delta W} = \underset{{\Delta W}}{\mathrm{argmin}} \;\; \mathcal{L}(Wx, \Delta Wx)
\]
where $\Delta W$ are the low-rank embedding matrices or prediction head, and $x$ is an input example.

We used a rank of 1024, compressing the model by 30\%. Then, we applied lightweight local distillation using our approach, with 64k randomly sampled Hausa sentences from InkubaMono\footnote{\url{https://huggingface.co/datasets/lelapa/Inkuba-Mono}} and the Aya-Dataset~\citep{singh2024aya}. We compared the models with MobiLLaMA~\citep{thawakar2024mobillama} and SmoLLM\footnote{\url{https://huggingface.co/blog/smollm}} using the AfriMMLU benchmark~\citep{afrimmlu}. As shown in Table~\ref{tab:inkuba-ch2}, the compressed InkubaLM model for Hausa retains 93\% of its base performance.

\begin{table}[h!]
    \centering
    {\small{
    \begin{tabular}{lcc}
    \toprule
    \textbf{Model} & \textbf{Model Size (Billion)} & \textbf{Accuracy} \\
    \midrule
    SmoLLM & 1.7 & 21.80 \\
    MobiLLaMA & 1.3 & 21.40 \\
    \midrule
    InkubaLM-base & 0.422 & 29.20 \\
    InkubaLM-30\% & 0.299 & 27.40 \\
    \bottomrule
    \end{tabular}
    }}
    \caption{Zero-shot performance on AfriMMLU for the 30\% compressed InkubaLM model for the Hausa language.}
    \label{tab:inkuba-ch2}
\end{table}

\section{Limitations}

\paragraph{Compatibility with quantization:}
Our pruning/compression method is theoretically complementary to quantization, since a model compressed with our approach can be further quantized. However, each step removes information, making it crucial to balance these approaches optimally. Addressing this challenge is essential for large-scale adoption in LLM libraries but is beyond the scope of this work and warrants a dedicated study.

\paragraph{Compromise between specificity and generality:}
Efficient LLM compression requires trade-offs between speed and generality. We prioritize generality over raw speed, avoiding GPU-specific kernels as used in semi-structured sparsity approaches. Efficiency also ties to carbon cost: our method works on CPUs as well as on
GPUs, enabling the reuse of older hardware and reducing reliance on GPUs with higher carbon costs.

\section{Conclusion}
We propose a new efficient compression method that
requires far less calibration data than most state-of-the-art approaches
and provides competitive performance for various models (dense and MoE Transformers, Mamba) and modalities (text, speech).
We show that combining SVD initialization with a joint Teacher and Student loss for local distillation enables fast convergence, and a bottom-up layer selection approach enables cost-effective compression.

%\addtocontents{toc}{\protect\clearpage} % <--- just debug stuff, ignore
\include{Chapters/Chapter03}
%\include{multiToC} % <--- just debug stuff, ignore for your documents
% ********************************************************************
% Backmatter
%*******************************************************
\ctparttext{%
\begin{flushright}{\slshape
    Le sens bénéfique de son [le caméléon] corps compressé, \\
    c'est s'appliquer à ne pas être encombrant. \\
    Au sens nocturne, c'est l'homme que l'on piétine. \\
    Le dos orné d'une crête symbolise la vigilance \\
    tel est son sens en diurne. \\
    En nocturne, c'est bomber la poitrine \\
    être avide de flatterie, plein d'orgueil et de boniments. \\
    La queue préhensile du caméléon est la prudence, \\
    une défense camouflée en un lieu imprévisible. \\
    En nocturne c'est un piège que le perfide \\
    recèle et qu'il traîne par devers lui} \\ \medskip
    --- Hampâté Bâ, Amadou. (1969). \\
    \mbox{\textit{Kaïdara : récit initiatique peul}.}
\end{flushright}}
\part{Activation Sparsity of Transformer Models}\label{pt:activation}
%*****************************************
\chapter{Cross-lingual Matryoshka Representation Learning across Speech and Text}\label{ch:mrl}
%*****************************************
\section{Introduction}
This chapter studies activation sparsity through the lens of Matryoshka Representation Learning
in embedding language models. We show that while Matryoshka representations offer a good trade-off
between efficiency and performance, they use only a subset of the available embedding space.
This leaves room to further improve how information is represented. In the next chapter, we
show that this observation also extends to autoregressive LLMs.

This chapter also reflects a desire to explore an engaging application beyond standard
English benchmarks. We ground the study in a concrete use case: retrieving French text
documents from Wolof speech queries. This task requires representations that bridge both
languages and modalities while remaining efficient. Matryoshka Representation Learning
allows us to explore this setting while continuing the thesis's focus on efficiency, by
training embeddings that can be used at different dimensions depending on the computational
budget.

Access to information is limited by two major barriers. First, a \emph{language barrier}: most online content is written in a few high-resource languages. Second, a \emph{modality barrier}: retrieval systems generally assume text queries, whereas many under-represented languages are primarily oral. Cascaded automatic speech recognition (ASR) and translation systems are costly and suffer from error propagation. We address both barriers with cross-lingual speech--text Matryoshka representations, enabling direct retrieval of text documents from speech queries with flexible accuracy--cost trade-offs.
We focus on Wolof and French as a socially grounded case study. Wolof is spoken primarily in Senegal and remains mainly oral. Because of the country's colonial history, much administrative, educational, and informational content relevant to Wolof speakers is available in French, while many Wolof speakers have limited French literacy. Information therefore exists as French text, but the natural query modality is Wolof speech.
We introduce the first bilingual speech--text Matryoshka embedding model for this setting and, to our knowledge, the first application of MRL to speech-to-text retrieval. We develop data-curation pipelines and new benchmarks, compare modeling architectures, and analyze how embedding dimension affects accuracy, storage, and latency. We also evaluate transfer learning to new tasks and study the effective rank of Matryoshka representations.

This study builds on Matryoshka Representation Learning, introduced in the state-of-the-art chapter, and its application to multilingual and cross-modal representations.
MRL~\citep{MRL} learns representations at multiple dimensions simultaneously, allowing the embedding dimension to be chosen at inference time. It has been generalized to vision--language and audio--visual models~\citep{hu2024matryoshkaquerytransformerlarge,cai2024matryoshkamultimodalmodels,cappellazzo2025adaptiveaudiovisualspeechrecognition}. We extend it to cross-lingual speech--text retrieval and analyze the performance, cost, and rank of its dimensions.

To achieve cross-modal speech and text representations, we draw inspiration from speech language models, which extend LLMs with speech understanding through a speech encoder~\citep{llamaspeech,pareras2025revisitingdirectspeechtotexttranslation,slm_methods}. We similarly add speech capabilities to embedding LLMs, an area underexplored relative to generative speech LLMs.

Cross-lingual text retrieval requires a strong multilingual speech representation model for Wolof and French.
XLM-R, multilingual BERT, and mBART~\citep{xlmr,mbert,mbart} demonstrate cross-lingual transfer but are not designed primarily for retrieval. 
NLLB-LLM2Vec~\citep{nllbllm2vec} produces cross-lingual representations for more than 200 languages, including Wolof and French.
We show that a model approximately ten times smaller, augmented with speech and trained mainly on supervised synthetic bilingual data, 
enables improved French and Wolof retrieval and transfers well to other tasks.

There are other approaches to achieving cross-modal representations,
such as CLIP~\citep{clip}, which uses modality-specific encoders with a contrastive objective. 
CLAP~\citep{clap} extends the design to audio and text.
We found that these approaches, called \textit{dual architectures}, work well for transcription retrieval but struggle with semantically demanding cross-lingual speech-to-document retrieval.

\section{Building a Wolof Speech Language Model}
\label{sec:wolof-speech-lm}

The retrieval system developed later in this chapter assumes a strong Wolof speech encoder. Such an encoder was not readily available: multilingual models cover Wolof only sparsely, and labeled Wolof speech is limited. We therefore first developed a Wolof-specific speech model in two stages. We continued self-supervised pretraining of HuBERT on a large collection of spontaneous Wolof speech, then connected the resulting encoder to a Wolof LLM. This section presents the complete development and evaluation of that model.

\subsection{Background and Motivation}

Self-supervised models such as wav2vec~\citep{wav2vec,wav2vec2} and HuBERT~\citep{hubert} learn general speech representations from unlabeled audio before supervised fine-tuning. wav2vec~2.0 uses a contrastive objective to predict discretized latent targets, requiring careful negative sampling. HuBERT instead clusters intermediate hidden representations and predicts their discrete cluster assignments with cross-entropy. Data2Vec~\citep{baevski2023efficientselfsupervisedlearningcontextualized} removes discretization and predicts masked continuous representations. We select HuBERT because its objective is simple, widely used, and supported by mature fine-tuning tools.

Multilingual models such as XLS-R and multilingual HuBERT~\citep{xlsr,mhubert} transfer representations across languages, and other work specifically targets African languages~\citep{gauthierSSL,afrihubert}. Performance on a particular language nevertheless depends on its share and diversity in the multilingual corpus. Wolof accounts for only about \(0.1\%\) of the African-centric pretraining data considered here. Because self-supervised learning does not require transcripts, scaling language-specific spontaneous audio is a practical alternative to relying on multilingual transfer alone.

\begin{figure}[ht]
    \centering
    \includegraphics[width=\linewidth]{papers/wolof_speech_lm/images/speech-pipeline.pdf}
    \caption[Wolof speech collection and filtering]{Collection of spontaneous Wolof speech for continued self-supervised pretraining. Raw web audio passes through source separation, speaker diarization, voice activity detection, and speech-quality filtering.}
    \label{fig:wolof-data-pipeline}
\end{figure}

\subsection{Wolof-Specific HuBERT Continued Pretraining}
\label{sec:wolof-hubert}

\subsubsection{Unsupervised Speech Data}

We manually curate 1.4\,TB of raw Wolof audio from public web sources, prioritizing natural, spontaneous speech by native speakers rather than read speech, which may contain hyper-articulation and artificial emotion. The filtering pipeline in Figure~\ref{fig:wolof-data-pipeline}, inspired by Emilia~\citep{emilia,emilialarge,amphion}, applies source separation, speaker diarization, and voice activity detection. We retain utterances lasting 3--30 seconds with DNSMOS above 3.2~\citep{dnsmos}. The result is 860 hours of intelligible, spontaneous Wolof speech.

\subsubsection{Initialization and Training}

We initialize the 90M-parameter HuBERT-base model~\citep{hubert}, originally pretrained on 960 hours of English LibriSpeech~\citep{librispeech}, and start with a fresh optimizer state because the original optimizer was unavailable. The Wolof corpus is nearly as large as the original English corpus, but initialization allows the model to reuse previously learned acoustic and phonetic structure.

Continued pretraining uses 32 A100 GPUs and the original HuBERT hyperparameters: 33 epochs, learning rate \(5\times10^{-4}\), and batches containing 87.5 seconds of audio. HuBERT-base was originally trained for 528 epochs; the shorter adaptation demonstrates the computational benefit of reusing a mature checkpoint rather than training from scratch.

\begin{insightbox}{Observation 1: Phonetic features transfer well across distant languages for continued pretraining}
Despite English and Wolof being typologically distant, English-pretrained HuBERT learns general phonetic features that transfer effectively to Wolof and require relatively few additional epochs.
\end{insightbox}

\subsection{Supervised ASR Fine-Tuning and Evaluation}
\label{sec:wolof-asr}

\subsubsection{Data and Fine-Tuning}

The supervised corpus combines FLEURS, ALFA~\citep{gauthierSSL}, Common Voice, Kallaama~\citep{kallama}, and Urban Bus~\citep{urbanbus}. We lowercase transcripts, remove punctuation, and convert digits to words. Official FLEURS splits are retained; the remaining corpora use random train/test partitions.

\begin{table}[ht]
    \centering
    \begin{tabular}{lrrrrrr}
        \toprule
        \textbf{Split} & \textbf{FLEURS} & \textbf{ALFA} & \textbf{CV} & \textbf{Kallaama} & \textbf{UB} & \textbf{Total} \\
        \midrule
        Train & 8.72 & 16.13 & 34.97 & 33.60 & 4.52 & \textbf{97.94} \\
        Test & 1.75 & 2.84 & 6.21 & 5.91 & 1.12 & \textbf{17.83} \\
        \bottomrule
    \end{tabular}
    \caption{Hours of supervised Wolof speech used for ASR fine-tuning and evaluation.}
    \label{tab:wolof-asr-data}
\end{table}

We convert the weights obtained through continued pretraining to the Hugging Face format and fine-tune with the standard CTC speech-recognition recipe. We compare three encoders fine-tuned on exactly the same supervised data: Meta/\allowbreak HuBERT-base, its version continually pretrained on Wolof, and Orange/\allowbreak HuBERT-base~\citep{afribert}, which was trained from scratch on 65,000 hours of African-language speech.

\subsubsection{Results and Interpretation}

\begin{table}[ht]
    \centering
    \begin{tabular}{lcc}
        \toprule
        \textbf{Model} & \textbf{Pretraining hours} & \textbf{WER \(\downarrow\)} \\
        \midrule
        Orange/HuBERT-base & 65,000 & 41.11 \\
        Meta/HuBERT-base & 960 & 39.48 \\
        Wolof HuBERT & \(960\rightarrow\mathbf{860}\) & \textbf{35.65} \\
        \bottomrule
    \end{tabular}
    \caption{ASR results after identical supervised fine-tuning. The arrow denotes English pretraining followed by Wolof continued pretraining.}
    \label{tab:wolof-hubert-wer}
\end{table}

Table~\ref{tab:wolof-hubert-wer} shows that language-specific continued pretraining reduces WER from 39.48 to 35.65 and outperforms the African-centric model despite using far less audio. Initialization transfers useful phonetic features, while the concentrated Wolof corpus is more valuable for Wolof than a much larger multilingual corpus in which Wolof is rare. Absolute WER remains relatively high partly because Wolof orthography is less standardized than that of high-resource written languages.

\begin{insightbox}{Observation 2: Do not rely on multilingualism alone; scale the monolingual corpus}
Orange/HuBERT-base underperforms on Wolof despite 65,000 hours of multilingual African speech containing Wolof and related languages. Cross-lingual transfer is useful, but scaling high-quality Wolof-specific data provides a stronger target-language model.
\end{insightbox}

\subsection{From a Speech Encoder to a Wolof Speech LLM}
\label{sec:wolof-slm}

HuBERT provides strong acoustic representations but its 90M parameters and CTC objective limit higher-level linguistic modeling. We connect it to a Wolof language model both to improve transcription and to enable speech translation.

\subsubsection{Training the Wolof Text LLM}

We adapt Qwen2.5-3B~\citep{qwen2,qwen2.5} using approximately two billion tokens. The mixture emphasizes bidirectional Wolof--English translation so that Wolof inputs can access knowledge learned during English pretraining, and includes high-quality monolingual Wolof text from Wolof-focused websites to improve generation.

\subsubsection{Linking HuBERT to the LLM}

\begin{figure}[ht]
    \centering
    \includegraphics[width=\linewidth]{papers/wolof_speech_lm/images/oolel_speech.pdf}
    \caption[Wolof Speech LLM architecture]{Wolof Speech LLM architecture. Features from every HuBERT layer are concatenated, projected by a lightweight aligner, combined with instruction-token embeddings, and passed to the LLM. Only the aligner is trained.}
    \label{fig:wolof-slm-architecture}
\end{figure}

We use late fusion~\citep{llava,llamaspeech}, illustrated in Figure~\ref{fig:wolof-slm-architecture}. Features from all HuBERT layers are concatenated and projected into the LLM embedding space by an MLP with one hidden layer with ReLU activation. The projected speech sequence is concatenated with text instruction embeddings and passed through the LLM. HuBERT and the LLM remain frozen; only the aligner is optimized. HuBERT already downsamples the waveform substantially, so no additional sequence downsampling is used.

Training examples use an instruction-following format in which the assistant transcribes or translates an audio input. The loss is computed only on assistant completions, with batch size 2 and learning rate \(10^{-4}\). The speech encoder is the CTC-fine-tuned Wolof HuBERT from Table~\ref{tab:wolof-hubert-wer}.

\subsubsection{Multi-step Speech Reasoning}

Besides direct transcription and translation, we test whether intermediate outputs improve performance. For ASR, the model may \textsc{phonemize} or \textsc{translate} before the \textsc{transcribe} step; phonemes provide an acoustic intermediate representation, whereas translation provides a semantic one. For translation, the model may first \textsc{transcribe} or \textsc{reformulate}. These variants test whether extra generation helps construct a stronger speech representation.

\subsubsection{Speech LLM Results}

We evaluate ASR with WER and character error rate (CER). Speech translation is evaluated on FLEURS using character \(n\)-gram F-score (ChRF) and BERTScore F1 (BS-F1), complementing surface overlap with semantic similarity.

\begin{table}[ht]
    \centering
    \begin{tabular}{lcc}
        \toprule
        \textbf{Generation order} & \textbf{WER \(\downarrow\)} & \textbf{CER \(\downarrow\)} \\
        \midrule
        \textsc{transcribe} & \textbf{29.09} & \textbf{15.26} \\
        \textsc{phonemize}\(\rightarrow\)\textsc{transcribe} & 34.05 & 19.08 \\
        \textsc{translate}\(\rightarrow\)\textsc{transcribe} & 29.48 & 15.95 \\
        \bottomrule
    \end{tabular}
    \caption{ASR results of the Wolof Speech LLM.}
    \label{tab:wolof-slm-asr}
\end{table}

\begin{table}[ht]
    \centering
    \begin{tabular}{lcc}
        \toprule
        \textbf{Generation order} & \textbf{ChRF \(\uparrow\)} & \textbf{BS-F1 \(\uparrow\)} \\
        \midrule
        \textsc{translate} & 33.08 & \textbf{79.78} \\
        \textsc{transcribe}\(\rightarrow\)\textsc{translate} & \textbf{33.79} & 79.73 \\
        \textsc{reformulate}\(\rightarrow\)\textsc{translate} & 33.59 & 77.54 \\
        \bottomrule
    \end{tabular}
    \caption{Speech-to-text translation results of the Wolof Speech LLM.}
    \label{tab:wolof-slm-translation}
\end{table}

Direct Speech LLM transcription obtains 29.09 WER, an 18.4\% relative improvement over the 35.65 WER of the HuBERT-only CTC model. Because HuBERT and the LLM are frozen while only the aligner is trained, the gain indicates that the LLM contributes useful Wolof linguistic knowledge rather than merely increasing acoustic capacity. The model also performs speech translation. Its translations are often semantic paraphrases rather than literal matches, explaining why BS-F1 is strong relative to ChRF.

Multi-step generation provides little benefit. Phonemization sometimes causes repetition loops in the 3B model, and transcription-before-translation yields only a small ChRF increase with essentially unchanged BS-F1. This agrees with evidence that the advantage of chain-of-thought diminishes as speech instruction data scales~\citep{pareras2025revisitingdirectspeechtotexttranslation}; larger language models may also be needed to exploit explicit intermediate reasoning.

\begin{insightbox}{Observation 3: Multi-step chain-of-thought does not improve this Speech LLM}
The LLM improves direct speech recognition and enables translation, but extra phonemization, transcription, translation, or reformulation steps provide little or no benefit. At this model and data scale, direct task supervision is the stronger operating point.
\end{insightbox}

\subsection{Implications for Cross-modal Retrieval}

The Wolof Speech LLM establishes the speech foundation used in the remainder of this chapter. Its development yields three practical lessons: collect substantial spontaneous monolingual audio; continue pretraining a strong self-supervised checkpoint instead of training from scratch; and connect a frozen LLM through a small trainable aligner when higher-level linguistic knowledge is needed. For retrieval, we retain the Wolof-adapted HuBERT encoder and late-fusion principle, but replace the generative LLM with a Matryoshka embedding LLM and train the speech projection for contrastive retrieval.

\section{Datasets}\label{sec:dataset-mrl}

Because Wolof speech--French text retrieval data are scarce, we rely mainly on synthetic documents. Here, a \emph{query} is any speech or text input used for retrieval, whether or not it is phrased as a question.

\begin{figure}[ht]
    \centering
    \includegraphics[width=\linewidth]{papers/MRL/images/pipeline.pdf}
    \caption[Speech data processing pipeline]{\textbf{Speech data pipeline.} Raw data are processed through source separation, diarization, voice activity detection, and quality filtering. Retained speech is transcribed before French stories, dialogues, and blog posts are generated.}
    \label{fig:pipeline-mrl}
\end{figure}

\subsection{Training Datasets}

\paragraph{Text data.}
We combine French mMARCO~\citep{mmarco}, whose queries are translated into Wolof; Senegalese French webpages, for which French queries are generated and translated; and French question-answering datasets with Wolof-translated questions. Wolof--French translation pairs encourage cross-lingual transfer. The result contains 1,176,908 query--document pairs and approximately 593 million document tokens.

\paragraph{Wolof speech queries.}
We reuse the 860-hour collection developed for Wolof HuBERT continued pretraining in Section~\ref{sec:wolof-hubert}. Figure~\ref{fig:pipeline-mrl} shows how this same filtered speech is extended into a retrieval-data pipeline: after transcription filtering, each query is paired with several kinds of French target documents.

\paragraph{Synthetic French documents.}
We transcribe queries using a Wolof speech language model~\citep{sy2025speechlanguagemodelsunderrepresented}. Perplexity and lexical-diversity filters retain roughly one quarter of the transcriptions. These are translated into French, and Gemini-2.5-Flash generates stories, dialogues, and web-style blog posts.

Synthetic targets are a practical way to scale supervision when parallel Wolof speech and French documents are unavailable. We considered reversing the pipeline by starting with natural French documents and synthesizing Wolof speech, but this would require a high-quality Wolof text-to-speech system, which is not available. Generating French documents is more reliable because current LLMs perform strongly in French. Using three document genres also reduces dependence on one synthetic style. Gemini's free API tier made generation at this scale affordable; preliminary trials indicated that open-source generators could also produce usable targets, but running them locally would have required substantially more GPU compute.

\paragraph{Instruction-following data.}
The same speech may be associated with a Wolof transcription, French translation, or French document. Instruction-following embedding models~\citep{peng-etal-2024-answer} resolve this ambiguity through a task description. We train one prompted model for document, translation, and transcription retrieval. The natural ASR corpora used for these additional tasks are detailed in Section~\ref{sec:wolof-asr}.

\subsection{Evaluation Datasets}

\begin{insightbox}{Evaluation artifacts: Wolof-to-French retrieval benchmarks}
\emph{Kallaama-Retrieval-Eval} and \emph{Fleurs-Retrieval-Eval} evaluate French document retrieval from Wolof speech or text.
\end{insightbox}

\paragraph{Fleurs-Retrieval-Eval.}
We begin with SIB-Fleurs~\citep{schmidt2025fleursslumassivelymultilingualbenchmark}, derived from FLEURS~\citep{fleurs}. French translations of Wolof speech are used to retrieve relevant French webpages, and failed or irrelevant results are manually removed. The 166 retained documents remain uncleaned to reflect deployment conditions.

\paragraph{Kallaama-Retrieval-Eval.}
FLEURS read speech is often hesitant and quiet, so we construct a benchmark from the natural conversations in Kallaama~\citep{kallama}. We select the 150 longest test utterances below 30 seconds, translate them, and use Gemini-2.5-Pro to generate a dialogue, blog post, and story for each. The result contains 150 speech and text queries, each paired with three French documents.

Both benchmarks are modest in size, which limits fine-grained statistical conclusions. Nevertheless, model inference and ranking are deterministic, the FLEURS targets are natural webpages collected from the web, and the Kallaama queries are spontaneous conversations representative of everyday spoken Wolof. We therefore use the two datasets as complementary evidence rather than treating either one as a comprehensive sample of the complete Wolof speech or French document distributions.

\section{Modeling}\label{sec:model-mrl}

\begin{figure}[ht]
    \centering
    \includegraphics[width=\linewidth]{papers/MRL/images/models.pdf}
    \caption{\textbf{Late-fusion and dual architectures.} Late fusion projects downsampled HuBERT features into the LLM embedding space and concatenates them with prompt embeddings. The dual architecture pools speech features and uses dimension-specific projections.}
    \label{fig:models-mrl}
\end{figure}

Given retrieval loss \(\mathcal{L}\) and dimensions \(\mathcal{M}=\{d_1,\ldots,d_m\}\), MRL optimizes
\begin{equation}
    \mathcal{L}_{\mathrm{MRL}}
    =\sum_{d\in\mathcal{M}}\mathcal{L}(\mathbf{Q}_{:d},\mathbf{D}_{:d}),
    \label{eq:mrl-loss}
\end{equation}
where \(\mathbf{Q}\) and \(\mathbf{D}\) are query and document embeddings and \({:d}\) slices the first \(d\) dimensions. We use InfoNCE~\citep{infonce} with in-batch negatives. The InfoNCE loss for a specific Matryoshka dimension is defined as:

\begin{equation}
    \mathcal{L}_{\mathrm{InfoNCE}}
    =\frac{\cos(q, d_{i})}{\sum_{d_{j}}\cos(q, d_{j})},
    \label{eq:infonce}
\end{equation}

\subsection{Text-only Embeddings}
We use Qwen3-0.6B-Embedding at dimensions 128, 256, 512, and 1024. It is first fine-tuned on text-only Wolof--French pairs with in-batch InfoNCE, teaching cross-lingual alignment before speech is introduced.

\subsection{Late Modality Fusion}
Late fusion~\citep{llava,llamaspeech} combines speech features and prompt-token embeddings before the LLM. We use Wolof-adapted HuBERT~\citep{sy2025speechlanguagemodelsunderrepresented}. Features from all 12 layers are concatenated, downsampled by a CNN, projected with \(\mathbf{W}\), concatenated with the prompt, and passed through the embedding LLM. Training uses Equation~\ref{eq:mrl-loss}; HuBERT and the LLM remain frozen while the CNN and projection are updated.

\begin{insightbox}{Model artifact: Cross-lingual speech and text representation models}
We introduce text-only and speech--text cross-lingual Matryoshka retrieval models for Wolof and French.
\end{insightbox}

\subsection{Dual Modality}
The dual architecture does not pass speech through the LLM. A learned query \(\mathbf{q}\) pools HuBERT features \(\mathbf{X}\):
\begin{equation}
    \operatorname{pool}(\mathbf{X})
    =\operatorname{softmax}\left(\frac{\mathbf{q}\mathbf{X}^{T}}{\sqrt{d}}\right)\mathbf{X}.
\end{equation}
Dimension-specific projections produce Matryoshka embeddings. This design is less expressive and cannot directly process prompts. \emph{Dual--Retrieval} uses InfoNCE while freezing the text model. \emph{Dual--Query Alignment} instead aligns speech with transcript representations using cosine-distance and \(\ell_1\) losses.

\section{Experiments}

\subsection{Setup}
We implement all four models with Sentence-Transformers. Each trains for one epoch with batch size 16, maximum length 2048, and learning rate \(3\times10^{-4}\). Although a single epoch may appear short, the dataset is large and both late-fusion and dual models plateau during that epoch; additional training produced no significant gains. The comparison therefore measures the architectures under the same compute budget rather than assuming that the dual models are fully optimized in an unconstrained setting. Fine-tuning Qwen3-0.6B uses eight A100 GPUs for approximately two hours, an estimated \$22.40 at \$1.40 per GPU-hour. The speech models train for approximately eight hours on one H100. They are cheaper because the large text model is frozen: late fusion updates only its small speech aligner, whereas the dual architecture updates the speech encoder and projection modules.

\subsection{Document Retrieval}
We evaluate normalized discounted cumulative gain (nDCG) at ranks 5 and 10. NLLB-LLM2Vec~\citep{nllbllm2vec} is an 8B-parameter text-only baseline with fixed 4096-dimensional embeddings. A pipelined baseline first transcribes speech and then uses our text-only model.

\begin{table}[ht]
    \centering
    \resizebox{\textwidth}{!}{%
    \begin{tabular}{lcccccccccc}
        \toprule
        \multirow{2}{*}{\textbf{Approach}} & \multicolumn{2}{c}{\textbf{4096}} & \multicolumn{2}{c}{\textbf{1024}} & \multicolumn{2}{c}{\textbf{512}} & \multicolumn{2}{c}{\textbf{256}} & \multicolumn{2}{c}{\textbf{128}} \\
        \cmidrule(lr){2-3}\cmidrule(lr){4-5}\cmidrule(lr){6-7}\cmidrule(lr){8-9}\cmidrule(lr){10-11}
        & @5 & @10 & @5 & @10 & @5 & @10 & @5 & @10 & @5 & @10 \\
        \midrule
        NLLB-LLM2Vec & 57.98 & 61.53 & -- & -- & -- & -- & -- & -- & -- & -- \\
        \rowcolor[HTML]{EEE0CC} Late-Fusion & -- & -- & \textbf{69.85} & \textbf{74.49} & 66.86 & 71.04 & 61.58 & 67.05 & 56.13 & 62.30 \\
        Pipelined & -- & -- & 57.09 & 62.82 & 54.07 & 59.14 & 45.60 & 51.56 & 41.21 & 47.34 \\
        Dual--Retrieval & -- & -- & 46.96 & 53.70 & 45.27 & 52.48 & 41.97 & 49.78 & 38.47 & 44.40 \\
        Dual--Alignment & -- & -- & 41.56 & 47.42 & 41.02 & 46.76 & 38.24 & 44.18 & 35.47 & 40.47 \\
        \bottomrule
    \end{tabular}}
    \caption[Kallaama retrieval results]{Document-retrieval nDCG on Kallaama-Retrieval-Eval.}
    \label{tab:kallaama-results-mrl}
\end{table}

\begin{table}[ht]
    \centering
    \resizebox{\textwidth}{!}{%
    \begin{tabular}{lcccccccccc}
        \toprule
        \multirow{2}{*}{\textbf{Approach}} & \multicolumn{2}{c}{\textbf{4096}} & \multicolumn{2}{c}{\textbf{1024}} & \multicolumn{2}{c}{\textbf{512}} & \multicolumn{2}{c}{\textbf{256}} & \multicolumn{2}{c}{\textbf{128}} \\
        \cmidrule(lr){2-3}\cmidrule(lr){4-5}\cmidrule(lr){6-7}\cmidrule(lr){8-9}\cmidrule(lr){10-11}
        & @5 & @10 & @5 & @10 & @5 & @10 & @5 & @10 & @5 & @10 \\
        \midrule
        NLLB-LLM2Vec & 55.98 & 59.43 & -- & -- & -- & -- & -- & -- & -- & -- \\
        \rowcolor[HTML]{EEE0CC} Late-Fusion & -- & -- & \textbf{57.89} & \textbf{61.19} & 55.03 & 58.90 & 50.87 & 54.43 & 41.56 & 46.60 \\
        Dual--Retrieval & -- & -- & 41.28 & 45.54 & 40.18 & 45.06 & 39.84 & 45.17 & 39.24 & 44.34 \\
        Dual--Alignment & -- & -- & 38.07 & 41.82 & 37.33 & 41.69 & 38.04 & 41.29 & 34.46 & 38.53 \\
        \bottomrule
    \end{tabular}}
    \caption{Document-retrieval nDCG on Fleurs-Retrieval-Eval.}
    \label{tab:fleurs-results-mrl}
\end{table}

Tables~\ref{tab:kallaama-results-mrl} and~\ref{tab:fleurs-results-mrl} show that late fusion performs best. On Kallaama, it outperforms NLLB-LLM2Vec despite being ten times smaller with embeddings four times shorter. It also outperforms the ASR pipeline. The dimension gap is larger on FLEURS, suggesting larger representations are more robust to low-quality speech.

The pipelined baseline already includes the first major source of cascading error: speech is transcribed before its text representation is retrieved. Adding machine translation would introduce another imperfect intermediate prediction, as well as additional latency, before retrieval. An ASR--translation--retrieval cascade should therefore be regarded as a less favorable operating point, especially for Wolof, where intermediate ASR and translation models are weaker than their high-resource counterparts. This expectation is an inference from the measured ASR--retrieval result rather than a directly measured three-stage baseline.

\begin{insightbox}{Observation 4: Late fusion outperforms dual architectures for cross-modal adaptation}
Late fusion provides stronger cross-modal generalization than dual encoders while training fewer parameters.
\end{insightbox}

\subsection{Speech Keyword Spotting}
We evaluate Urban Bus~\citep{urbanbus}, retrieving the station or landmark spoken by the user. This is equivalent to transcription retrieval.

\begin{table}[ht]
    \centering
    \begin{tabular}{llcccc}
        \toprule
        \multicolumn{2}{c}{\textbf{Approach}} & \textbf{1024} & \textbf{512} & \textbf{256} & \textbf{128} \\
        \midrule
        \cellcolor[HTML]{EEE0CC} & \cellcolor[HTML]{EEE0CC}F1 & \cellcolor[HTML]{EEE0CC}\textbf{88.79} & \cellcolor[HTML]{EEE0CC}\textbf{84.85} & \cellcolor[HTML]{EEE0CC}\textbf{79.80} & \cellcolor[HTML]{EEE0CC}\textbf{76.86} \\
        \multirow{-2}{*}{\cellcolor[HTML]{EEE0CC}Late-Fusion} & \cellcolor[HTML]{EEE0CC}Recall & \cellcolor[HTML]{EEE0CC}\textbf{89.79} & \cellcolor[HTML]{EEE0CC}\textbf{86.49} & \cellcolor[HTML]{EEE0CC}\textbf{81.98} & \cellcolor[HTML]{EEE0CC}\textbf{79.88} \\
        \multirow{2}{*}{Dual--Retrieval} & F1 & 68.64 & 58.33 & 50.86 & 44.59 \\
        & Recall & 71.17 & 62.76 & 56.76 & 52.25 \\
        \multirow{2}{*}{Dual--Alignment} & F1 & 76.07 & 67.67 & 67.67 & 65.22 \\
        & Recall & 77.78 & 70.87 & 70.87 & 68.47 \\
        \bottomrule
    \end{tabular}
    \caption{Keyword-spotting performance across embedding dimensions.}
    \label{tab:kws-mrl}
\end{table}

Late fusion performs best at every dimension. Dual encoders nevertheless perform much better here than on document retrieval, and query alignment outperforms direct retrieval training.

\begin{insightbox}{Observation 5: Dual models remain useful for less semantically demanding tasks}
Dual encoders remain viable for transcription retrieval, which requires less semantic understanding than document retrieval.
\end{insightbox}

\subsection{Generalization to an Unseen Task}\label{sec:unseen-task-mrl}
We evaluate speech intent detection on WolBanking77~\citep{wolbanking}, which has ten banking intents. Zero-shot classification selects the label most similar to a speech request. In the \(n\)-shot setting, we follow SetFit~\citep{setfit}: contrastive fine-tuning improves representations before a logistic-regression head is trained.

\begin{table}[ht]
    \centering
    \begin{tabular}{ccc}
        \toprule
        \textbf{\(n\)-shot} & \textbf{F1} & \textbf{Recall} \\
        \midrule
        0 & 44.79 & 50.64 \\ 1 & 56.29 & 58.50 \\ 2 & 60.40 & 61.93 \\
        4 & 87.18 & 88.46 \\ 8 & 92.54 & 92.51 \\ 16 & 96.11 & 96.10 \\
        \bottomrule
    \end{tabular}
    \caption{Few-shot speech intent detection at dimension 1024.}
    \label{tab:fewshot-mrl}
\end{table}

Zero-shot F1 is substantially above the 10\% random baseline and reaches 96.11 with 16 examples per class. Figure~\ref{fig:sid-dimensions-mrl} shows that larger embeddings adapt faster, while smaller ones catch up as the amount of data increases.

\begin{insightbox}{Observation 6: Larger embeddings adapt quickly; smaller embeddings need more data}
Larger embeddings adapt rapidly to new tasks; smaller embeddings require more examples but eventually approach their performance.
\end{insightbox}

\section{Analysis}\label{sec:analysis-mrl}

\begin{figure}[ht]
    \centering
    \includegraphics[width=0.8\linewidth]{papers/MRL/images/energy_evolution.pdf}
    \caption[Dimensions needed for an energy ratio]{Dimensions needed for a given energy ratio.}
    \label{fig:energy-evolution-mrl}
\end{figure}

\begin{figure}[ht]
    \centering
    \includegraphics[width=0.8\linewidth]{papers/MRL/images/energy_100p.pdf}
    \caption{Dimensions needed for full energy.}
    \label{fig:energy-full-mrl}
\end{figure}

\begin{figure}[ht]
    \centering
    \includegraphics[width=0.8\linewidth]{papers/MRL/images/recall_vs_nshot.pdf}
    \caption{Recall by embedding dimension and number of examples.}
    \label{fig:sid-dimensions-mrl}
\end{figure}

\subsection{Speech Quality in FLEURS}
FLEURS asks speakers to read translated texts, but Wolof's standardized writing system is not widely taught. Its hesitations thus primarily reflect difficulty reading Wolof rather than the discourse hesitations that naturally occur in spontaneous conversation. Compared with Kallaama, FLEURS contains more hesitations and lower volume. It averages 7.51 characters per second versus 16.88 for Kallaama, while mean loudness is \(-49.01\) dB versus \(-23.81\) dB. This helps explain Table~\ref{tab:fleurs-results-mrl} and the benefit of larger dimensions for low-quality speech.

\begin{insightbox}{Observation 7: Wolof speech in FLEURS is unnatural}
Read Wolof in FLEURS is less natural, more hesitant, and quieter than spontaneous Kallaama speech.
\end{insightbox}

\subsection{Prompting}
Without prompts, late fusion obtains 67.07/71.68 nDCG@5/10 on Kallaama at dimension 1024, compared with 46.96/53.70 for the dual model. Its advantage is therefore not explained only by prompt support. Task-aligned prompts still help: the document prompt produces the best document scores (68.85/74.49), while the transcription prompt produces the best keyword scores (88.79 F1, 89.79 recall).

\subsection{Rank of Matryoshka Dimensions}
We compute the cumulative energy ratio
\begin{equation}
    R(k)=\frac{\sum_{i=1}^{k}\lambda_i}{\sum_{j=1}^{d}\lambda_j},
\end{equation}
where the \(\lambda_i\) are descending eigenvalues of the representation covariance matrix on Kallaama. Figure~\ref{fig:energy-evolution-mrl} shows that larger embeddings reach full energy using a smaller fraction of coordinates and have lower relative rank. Yet smaller prefix embeddings do not match their accuracy, suggesting slicing discards critical information. Prior work similarly identifies gradient variance and rigid slicing as MRL limitations~\citep{zhang-etal-2025-smec,wen2025matryoshkarevisitingsparsecoding}. Figure~\ref{fig:energy-full-mrl} also shows that documents have higher ranks than speech queries; dimensions 128 and 256 are full-rank.

\begin{insightbox}{Observation 8: Information is distributed unevenly across Matryoshka dimensions}
Large representations concentrate information in relatively few components, leaving scope for better compression than prefix slicing.
\end{insightbox}

\subsection{Deployment Costs}
\begin{figure}[ht]
    \centering
    \includegraphics[width=\linewidth]{papers/MRL/images/costs.pdf}
    \caption{Indexing, storage, query-time, and accuracy trade-offs across dimensions.}
    \label{fig:costs-mrl}
\end{figure}
We index float16 embeddings with \texttt{chromadb} and measure indexing throughput, disk usage, and single-result query latency. Figure~\ref{fig:costs-mrl} shows diminishing returns: larger embeddings provide smaller accuracy gains while increasing storage. Smaller embeddings accelerate retrieval substantially, although indexing speed improves more gradually. The best dimension depends on the task and compute budget.

The stronger degradation of small dimensions on FLEURS does not imply that low-dimensional embeddings work only with clean or near-perfect recordings. It is the standard efficiency--accuracy trade-off becoming steeper under more challenging acoustic conditions: larger vectors preserve more information and are more robust, while smaller vectors retain their storage and speed advantages at the cost of a larger accuracy drop. Deployment should therefore select the embedding dimension in conjunction with expected input quality rather than treating one dimension as universally optimal.

\section{Limitations}

\paragraph{Synthetic data.}
Most documents are generated from transcribed and translated speech. This enables scale but may not reproduce real-world diversity and noise. It also assumes usable ASR and translation systems, which many languages lack.

\paragraph{Beyond Matryoshka slicing.}
Only a small subset of coordinates represents most information, suggesting prefix slicing is suboptimal. Sparse-coding alternatives~\citep{wen2025matryoshkarevisitingsparsecoding} or dynamic structured sparsity may offer more flexible compression.

\section{Conclusion}
We introduced cross-lingual speech--text Matryoshka embedding models that retrieve French documents directly from Wolof speech. Late modality fusion consistently outperforms dual encoders and an ASR pipeline, while the learned representations transfer to keyword spotting and intent detection.

The dimension and rank analyses expose a central limitation: information is concentrated in a small number of components, yet prefix-truncated embeddings do not preserve enough of it to match the largest representation. This connects adaptive representation learning to activation sparsity and motivates more flexible mechanisms for selecting informative dimensions at inference time.

\include{Chapters/Chapter05}
\ctparttext{%
\begin{flushright}{\slshape
    « Ce que tu viens d'apprendre, retiens-le \\
    toi qui sais que la science par essence \\
    vaut plus que l'ambre et le corail \\
    et plus cher que l'or pur. \\
    Tu as longtemps médité et longuement cherché. \\
    Qui s'applique à chercher finit par trouver. \\
    Tu n'avais vu que les signes du phénomène; \\
    et maintenant tu as le sens du premier secret \\
    au pays des génies-nains qui appartient à Kaydara \\
    le lointain, le bien proche Kaydara ».} \\ \medskip
    --- Hampâté Bâ, Amadou. (1969). \\
    \mbox{\textit{Kaïdara : récit initiatique peul}.}
\end{flushright}}
\part{Conclusion and Future Work}\label{pt:conclusion}
%********************************************************************
% Overall conclusion and future work
%********************************************************************
\chapter{Overall Conclusion and Future Work}

\section{Overall Conclusion}
The promises of scaling laws, where scaling up model size gives rise to new capabilities, 
are colliding with a harsh reality: training these large models and running inference
are costly. The central research question is how to scale language models efficiently,
both for training and inference. It is fortunate that pre-trained LLMs exhibit
redundancy at different levels that can be exploited to improve their efficiency
during post-training or pretraining. One example of
such redundancy is the low-rank structure of pretrained Transformer activations.
Our contributions exploit this observation to improve LLM efficiency in two main
directions: \textbf{model compression} using activations and \textbf{activation sparsity}.

For post-training \textbf{model compression}, we propose \textbf{Lillama}, a method that locally distills the
activations of Transformer layers into low-rank weights. This local optimization makes
the approach more efficient than prior work: it enables compressing, on a single GPU,
LLMs that would not otherwise fit on one. To achieve speedup, the method does not require
specialized GPU kernels, unlike weight-sparsity approaches such as SparseGPT~\citep{frantar2023sparsegpt} or Wanda~\citep{sun2024wanda}.
The approach is also sample-efficient, requiring only 13M tokens, compared to methods
like Sheared-Llama \citep{xia2024sheared}, or Minitron \citep{muralidharan2024compact,
sreenivas2024minitron}, which use over 50B tokens for compression. Sample efficiency is an important property in post-training compression, as it
directly affects the fairness of the approach---for instance, its viability for
under-represented languages. This is why we specifically studied Whisper compression in
a low-resource setting, using Bambara as a case study. In this work, \textbf{BaldWhisper}, we
again distill input/output embedding activations into low-rank weights, but push
compression further by reducing the number of layers. To preserve the method's sample
efficiency, we do this by merging layers rather than removing them,
to stay close to the original model.
The resulting model is twice as fast while maintaining most of the performance.

For \textbf{activation sparsity}, we first study the eigenstructure of \textbf{Matryoshka
embedding LLMs} in multilingual and multimodal representation learning. We show
that Matryoshka embeddings, viewed as a form of structured activation sparsity,
leave room for more effective use of the representation space: most of the
activation energy is concentrated in relatively few principal components.
The fixed slicing operator applied to the embedding vectors
does not preserve all the information needed for downstream tasks while
introducing redundancy. We make the same observation for dense MLPs in
autoregressive LLMs: their intermediate activations occupy a subspace
substantially smaller than the allocated dimension. Sparse MoEs, however,
project each token into smaller subspaces of the available intermediate space,
and we show that these smaller subspaces are used more fully.
These observations motivate an approach where the MLP has a large intermediate
dimension but each token adaptively selects only the intermediate coordinates it needs
rather than activating the whole available space.
We introduce \textbf{ZeroNeuron (ZoN)}, an input-dependent
gating mechanism that learns which coordinates to activate while constraining the average active dimension.
Compared to fine-grained sparse mixture-of-experts and sparse memory layers,
ZoN gets rid of TopK and predicts sparsity pointwise, which
makes it not only more compute-adaptive but also more communication-efficient
in distributed training, compared to TopK-based sparse mixture-of-experts and
sparse memory layers that require an all-to-all token dispatch and all-to-all return.
With approximately 600M active parameters per
token, the 1.2B ZoN model matches the aggregate performance of its dense
counterpart and outperforms a dense model of a comparable number of active parameters.
These results show that model capacity can be preserved while reducing
the parameters activated for each input, and that such sparsity
can be exploited by GPU kernels to accelerate training and inference.

\section{Future Work}
We believe in sparsity, and we expect the most performant machine learning systems
of the future to be sparsely activated. The remaining question is how
this sparsity should be designed. We do not believe that TopK-based
mixtures of experts or sparse memory layers are the right long-term
answer: there is no reason to impose the same activation budget on
every token, regardless of the task and its difficulty. Moreover, under
current hardware constraints, distributed TopK routing incurs substantial
cross-GPU expert communication. This is one of the biggest challenges in
scaling fine-grained sparsity, and approaches such as LatentMoE
~\citep{elango2026latentmoeoptimalaccuracyflop} are trying to address
it. Maybe this will no longer be a problem in a few years, if the next generation
of hardware offers individual GPUs with terabytes of VRAM and the capacity to train
LLMs with trillions of parameters on a single device.
But beyond the engineering issues, TopK-based
routing also faces the stability problems discussed earlier, including
dead experts and expert collapse. Rather than putting patches everywhere
to keep TopK-based sparsity working, we believe it can be beneficial to fundamentally
revisit how sparsity is learned. We believe that such approaches must be adaptive, communication-efficient,
and stable by design. ZoN is our first step in this direction. We plan
to investigate its scalability through larger models and longer training
runs, while developing GPU kernels that translate its learned sparsity
into actual computational savings.
%********************************************************************
% Other Stuff in the Back
%*******************************************************
\pdfbookmark[-1]{Back Matter}{back} % Dummy part bookmark for back matter
\cleardoublepage\include{FrontBackmatter/Bibliography}
% ********************************************************************
% Game Over: Restore, Restart, or Quit?
%*******************************************************
\end{document}
% ********************************************************************
